\documentclass[11pt,a4paper]{article}
\usepackage[utf8]{inputenc}
\usepackage[margin=1in]{geometry}
\usepackage{graphicx}
\usepackage{longtable}
\usepackage{booktabs}
\usepackage{hyperref}
\usepackage{array}
\usepackage{fancyhdr}
\usepackage{amsmath}
\usepackage{titlesec}
\usepackage{color}
\usepackage{float}
\usepackage{caption}
\usepackage{microtype}

\pagestyle{fancy}
\fancyhf{}
\fancyhead[L]{OzFlux Tower Network: Detailed Site Reports}
\fancyhead[R]{\thepage}
\fancyfoot[C]{Confidential - Research Use Only}
\renewcommand{\headrulewidth}{0.4pt}
\renewcommand{\footrulewidth}{0pt}

\begin{document}

\title{\textbf{OzFlux Tower Network: Detailed Site Reports (V2)}}
\author{OzFlux Team}
\date{\today}
\maketitle

\begin{abstract}
OzFlux is a national network of flux towers that provides continuous measurements of carbon, water, and energy exchanges between terrestrial ecosystems and the atmosphere across Australia and New Zealand. This detailed report merges physical configuration diagrams, scientific purposes, site descriptions, instrument tables, and on-site photographs to document each monitoring station.
\end{abstract}

Total sites documented: 31

\newpage

\tableofcontents
\newpage

\section{Site Profile: Adelaide River}

\subsection*{Physical Configuration \& Soil Sensor Depths}

\begin{figure}[H]
\centering
  \includegraphics[width=0.85\textwidth,height=0.45\textheight,keepaspectratio]{/home/sanjays/et97_scratch2/oldscratch/Ozflux_data_full/L6/soil_depth_plots/FluxTower_Diagrams/AdelaideRiver/AdelaideRiver_diagram.png}
  \caption*{Soil sensor depth profile and tower instrument heights.}
\end{figure}

\newpage

\subsection*{Site Description}

\textbf{Site Description}

The Adelaide River flux station was located approximately 10.5km south east of Bachelor, Northern    Territory (GPS coordinates: -13.0769, 131.1178).

The flux tower site was classified as Savanna  dominated by \textit{Eucalyptus tectifica}and \textit{Planchonia  careya.}

Elevation of the site was close to 90m and mean annual  precipitation at a nearby Bureau of Meteorology site is 1730mm. Maximum  temperatures range from 31.4$^\circ$C (in June) to 36.8$^\circ$C (in  October) while minimum temperatures range from 16.2$^\circ$C (in July) to  25.1$^\circ$C (in December). Maximum temperature vary seasonally by  approximately 5.4$^\circ$C and minimum temperatures vary by approximately  8.9$^\circ$C.

The instrument mast was 15 meters  tall. Heat, water vapour and carbon dioxide measurements were taken using the  open-path eddy flux technique. Temperature, humidity, wind speed, wind  direction, rainfall, incoming and reflected shortwave radiation and net  radiation were measured above the canopy. Soil heat fluxes are measured and soil  moisture content was gathered using time domain reflectometry.

\subsection*{Purpose \& Scientific Goals}

\textbf{Purpose}

The  primary purpose of the Adelaide River Flux Station was to:

\begin{itemize}
  \item Provide  information as part of a larger network of flux towers established along the  North Australian Tropical Transect (NATT) gradient, which extends \textasciitilde{}1000km  south from Darwin 12.5 S.  The towers provide flux data across the savanna’s heterogeneous ‘top end’  ecosystems including open forest savanna, open savanna woodland and seasonally  inundated floodplains.
  \item Examine spatial patterns and processes of  land-surface-atmosphere exchanges (radiation, heat, moisture, CO2 and other  trace gasses) across scales from leaf to landscape scales within Australian  savannas.
  \item Determine  the climate and ecosystem characteristics (physical structure,  species composition, physiological function) that drive spatial and temporal variations of  carbon, water and energy fluxes from north Australian savanna.
  \item Determine if fluxes of carbon, water vapour and  heat over the various ecosystems as derived from the various measurement  techniques can be combined to form a comprehensive and consistent estimate of  the regional fluxes and budgets across the landscape.
  \item Provide  longer term measurements for future projects.
\end{itemize}

\subsection*{Site Photographs}

\begin{figure}[H]
\centering
  \includegraphics[width=0.8\textwidth,height=0.4\textheight,keepaspectratio]{/home/sanjays/et97_scratch2/oldscratch/Ozflux_data_full/ozflux_mirror/www.ozflux.org.au/monitoringsites/adelaideriver/images/AR2.JPG}
  \caption*{Photo 1: The Adelaide River Flux Tower}
\end{figure}

\begin{figure}[H]
\centering
  \includegraphics[width=0.8\textwidth,height=0.4\textheight,keepaspectratio]{/home/sanjays/et97_scratch2/oldscratch/Ozflux_data_full/ozflux_mirror/www.ozflux.org.au/monitoringsites/adelaideriver/images/AR1.jpg}
  \caption*{Photo 2: The  understory of the tower site}
\end{figure}

\begin{figure}[H]
\centering
  \includegraphics[width=0.8\textwidth,height=0.4\textheight,keepaspectratio]{/home/sanjays/et97_scratch2/oldscratch/Ozflux_data_full/ozflux_mirror/www.ozflux.org.au/monitoringsites/adelaideriver/images/AR3.jpg}
  \caption*{Photo 3: Tower  Instrumentation}
\end{figure}

\subsection*{Measurement Instrumentation}

\begin{longtable}{p{4.5cm}p{3.5cm}p{3.5cm}p{3.5cm}}
\toprule
\textbf{Instrument Type} & \textbf{Make} & \textbf{Model} & \textbf{Situation} \\
\midrule
\endhead
\bottomrule
\endfoot
Data Logger & Campbell Scientific & CR3000 & Ground \\
Open path  CO2 and H2O & LiCor & Li-7500 & 15m \\
Net Radiometer & Kipp \& Zonen & CNR1 & 15m \\
Rain  Gauge & Campbell Scientific & CS702 & Ground \\
Soil  Heat Flux (two replicates) & Campbell Scientific & HFT3 & Ground, -0.15m \\
Soil  moisture (two replicates) & Soil  moisture (two replicates) & CS616 & -0.05,  -0.50m \\
Sonic anemometer                 - wind velocities (u,v,w)                  - sonic   temperature & Campbell  Scientific & CSAT3 & 15m \\
temperature and Relative Humidity & Vaisalla & HMP45A & 15m \\
\end{longtable}


\subsection*{Tower Global Metadata}

\begin{longtable}{p{4.5cm}p{10.5cm}}
\toprule
\textbf{Attribute} & \textbf{Value} \\
\midrule
\endhead
\bottomrule
\endfoot
\textbf{PI} & Jason Beringer and Lindsay Hutley \\
\textbf{Institution} & University of Western Australia \\
\textbf{Latitude} & -13.0769 \\
\textbf{Longitude} & 131.1178 \\
\textbf{Elevation} & N/A \\
\textbf{Soil Type} & red kandasol \\
\textbf{Vegetation} & Woody Savanna \\
\textbf{Temporal Coverage} & 2007-10-17 11:30:00 to 2009-05-24 06:00:00 \\
\textbf{Data Link} & \url{http://data.ozflux.org.au/} \\
\end{longtable}

\newpage

\section{Site Profile: Alpine Peatland}

\subsection*{Physical Configuration \& Soil Sensor Depths}

\begin{figure}[H]
\centering
  \includegraphics[width=0.85\textwidth,height=0.45\textheight,keepaspectratio]{/home/sanjays/et97_scratch2/oldscratch/Ozflux_data_full/L6/soil_depth_plots/FluxTower_Diagrams/AlpinePeatland/AlpinePeatland_diagram.png}
  \caption*{Soil sensor depth profile and tower instrument heights.}
\end{figure}

\newpage

\subsection*{Site Description}

\textbf{Site Description}

The  Alpine Peatland flux station is located on the Bogong High Plains near Falls  Creek in the Victorian Alpine National Parks.

The  vegetation community is Alpine \textit{Sphagnum} peatland, also known as Closed Heath, composed of Sphagnum moss\textit{(Sphagnum cristatum)}in association with candle heath\textit{(Richea continentis),}alpine baeckea\textit{(Baeckea gunniana)}and rope rush\textit{(Empodisma minus).} The underlying soil is peat.

Elevation of the site is 1670m and mean annual precipitation is 1274mm.  Mean annual average minimum and maximum temperatures are 2.6$^\circ$C and 9.4$^\circ$C and  the area is snow covered for an average of 3 months per year.

The Victorian  Alpine National Park is managed for conservation, tourism and as a water  catchment for hydroelectricity production. Historic disturbance includes  grazing, fire and infrastructure development.

\subsection*{Purpose \& Scientific Goals}

\textbf{Purpose}

The purpose of the Alpine Peatland flux station is:

\begin{itemize}
  \item to directly measure evapotranspiration, supporting detailed long term hydrological measurements that are ongoing at this site (Karis et al., 2016; Silvester, 2006; Silvester, 2009).
  \item to investigate carbon and energy fluxes
  \item to add to the range of ecosystems encompassed by the OzFlux network as there has been previously little coverage of alpine, snow-covered landscapes or wetlands.
\end{itemize}

\subsection*{Site Photographs}

\begin{figure}[H]
\centering
  \includegraphics[width=0.8\textwidth,height=0.4\textheight,keepaspectratio]{/home/sanjays/et97_scratch2/oldscratch/Ozflux_data_full/ozflux_mirror/www.ozflux.org.au/monitoringsites/alpinepeatland/images/teamWithTower.jpg}
  \caption*{Photo 1: The Alpine Peatland Flux tower and its team.}
\end{figure}

\begin{figure}[H]
\centering
  \includegraphics[width=0.8\textwidth,height=0.4\textheight,keepaspectratio]{/home/sanjays/et97_scratch2/oldscratch/Ozflux_data_full/ozflux_mirror/www.ozflux.org.au/monitoringsites/alpinepeatland/images/fluxInstruments.jpg}
  \caption*{Photo 2: Flux instrumetation.}
\end{figure}

\begin{figure}[H]
\centering
  \includegraphics[width=0.8\textwidth,height=0.4\textheight,keepaspectratio]{/home/sanjays/et97_scratch2/oldscratch/Ozflux_data_full/ozflux_mirror/www.ozflux.org.au/monitoringsites/alpinepeatland/images/towerInSnow_July2017.jpg}
  \caption*{Photo 3: The tower in the snow.}
\end{figure}

\subsection*{Measurement Instrumentation}

\begin{longtable}{p{4.5cm}p{3.5cm}p{3.5cm}p{3.5cm}}
\toprule
\textbf{Instrument Type} & \textbf{Make} & \textbf{Model} & \textbf{location} \\
\midrule
\endhead
\bottomrule
\endfoot
Open path  CO2, H2O & Campbell Scientific & IRGASON & 2.2m \\
Sonic anemometer & Campbell Scientific & IRGASON & 2.2m \\
Net Radiometer & Kipp and Zonen & CNR1 & 2.2m \\
Rh+T & Vaisala & HMP60 & 2.2m \\
Atmospheric Pressure & Vaisala & PTB110 & 2.2m \\
Anemometer & Vaisala & WAA151 & 2.2m \\
Wind Vane & Vector Instruments & W200P & 2.2m \\
\end{longtable}


\subsection*{Tower Global Metadata}

\begin{longtable}{p{4.5cm}p{10.5cm}}
\toprule
\textbf{Attribute} & \textbf{Value} \\
\midrule
\endhead
\bottomrule
\endfoot
\textbf{PI} & Samantha Grover \\
\textbf{Institution} &  \\
\textbf{Latitude} & -36.862222 \\
\textbf{Longitude} & 147.320833 \\
\textbf{Elevation} & N/A \\
\textbf{Soil Type} &  \\
\textbf{Vegetation} &  \\
\textbf{Temporal Coverage} & 2017-04-12 18:30:00 to 2022-06-20 23:30:00 \\
\textbf{Data Link} & \url{http://data.ozflux.org.au/} \\
\end{longtable}

\newpage

\section{Site Profile: Boyagin}

\subsection*{Physical Configuration \& Soil Sensor Depths}

\begin{figure}[H]
\centering
  \includegraphics[width=0.85\textwidth,height=0.45\textheight,keepaspectratio]{/home/sanjays/et97_scratch2/oldscratch/Ozflux_data_full/L6/soil_depth_plots/FluxTower_Diagrams/Boyagin/Boyagin_diagram.png}
  \caption*{Soil sensor depth profile and tower instrument heights.}
\end{figure}

\newpage

\subsection*{Site Description}

\textbf{Site Description}

The Boyagin Wandoo Woodland flux station is located  in the wheatbelt of Western Australia about 200km southeast from Perth,
            Western Australia (GPS coordinates: -32.477093, 116.938559; elevation 484m).

The flux tower site is located within an area of wandoo woodland. The surrounding area is dominated by broadacre farming practices. Elevation of the site is close to 484m and mean annual precipitation at a nearby Bureau of Meteorology site at Pingelly measures 446mm.

Maximum temperatures range from 15.3$^\circ$C (in July) to 31.9$^\circ$C (in Jan), while minimum temperatures range from 5.5$^\circ$C (in July) to 15.5$^\circ$C (in Jan).

The instrument mast is 4 meters tall. Heat, water vapour, and carbon dioxide measurements are taken using the open-path eddy flux technique. Temperature, humidity, wind speed, wind direction, rainfall and net radiation are measured. Soil heat fluxes are measured and soil moisture content are collected.

\subsection*{Purpose \& Scientific Goals}

\textbf{Purpose}

The purpose of the Boyagin Wandoo Woodland Flux Station is to:

\begin{itemize}
  \item Monitor and determine the balance of environmental demands for water yields, agricultural productivity, GHG budgets and biodiversity within a catchment landscape.
  \item Provide information to establish a modelling tool for GHG and water fluxes across various land use types, in order to benefit land management practices in the wheatbelt of Western Australia.
\end{itemize}

The tower is a crucial component of the UWA critical zone observatory which will be the site of a multidisciplinary approach to understanding the landscape dynamics.

\subsection*{Site Photographs}

\begin{figure}[H]
\centering
  \includegraphics[width=0.8\textwidth,height=0.4\textheight,keepaspectratio]{/home/sanjays/et97_scratch2/oldscratch/Ozflux_data_full/ozflux_mirror/www.ozflux.org.au/monitoringsites/boyagin/images/IMG_1335.jpg}
  \caption*{Photo 1: The Boyagin flux tower team, Atbin, Ian, Tim, Matt, and Jason}
\end{figure}

\begin{figure}[H]
\centering
  \includegraphics[width=0.8\textwidth,height=0.4\textheight,keepaspectratio]{/home/sanjays/et97_scratch2/oldscratch/Ozflux_data_full/ozflux_mirror/www.ozflux.org.au/monitoringsites/boyagin/images/DJI_0036.jpg}
  \caption*{Photo 2: Start of the instrument installation}
\end{figure}

\begin{figure}[H]
\centering
  \includegraphics[width=0.8\textwidth,height=0.4\textheight,keepaspectratio]{/home/sanjays/et97_scratch2/oldscratch/Ozflux_data_full/ozflux_mirror/www.ozflux.org.au/monitoringsites/boyagin/images/IMG_1375.jpg}
  \caption*{Photo 3: Instrumentation at the top of the tower}
\end{figure}

\begin{figure}[H]
\centering
  \includegraphics[width=0.8\textwidth,height=0.4\textheight,keepaspectratio]{/home/sanjays/et97_scratch2/oldscratch/Ozflux_data_full/ozflux_mirror/www.ozflux.org.au/monitoringsites/boyagin/images/IMG_1372.jpg}
  \caption*{Photo 4: Completed tower...}
\end{figure}

\begin{figure}[H]
\centering
  \includegraphics[width=0.8\textwidth,height=0.4\textheight,keepaspectratio]{/home/sanjays/et97_scratch2/oldscratch/Ozflux_data_full/ozflux_mirror/www.ozflux.org.au/monitoringsites/boyagin/images/IMG_1327.jpg}
  \caption*{Photo 5: ..and as the sun sets}
\end{figure}

\subsection*{Measurement Instrumentation}

\begin{longtable}{p{4.5cm}p{3.5cm}p{3.5cm}p{3.5cm}}
\toprule
\textbf{Instrument Type} & \textbf{Make} & \textbf{Model} & \textbf{Situation} \\
\midrule
\endhead
\bottomrule
\endfoot
Atmospheric Pressure & Vaisala & PTB110 & 2.5m \\
Net Radiation & Kipp and Zonen & CNR4 & 4m \\
Data logger & Campbell Scientific & CR3000 & Ground \\
Open Path CO2 \& H2O & LI-COR & LI-7500 & 2.5m \\
Rain gauge & Hydrological Services & TB5 & Ground \\
Soil heat flux (two replicates) & Hukseflux self calibrating probes & HFP01 & 0.1m \\
Soil temperature & Campbell Scientific & TCAV-L & 0.08m \\
Soil moisture & Campbell Scientific & CS616 & 0.1, 0.2, 0.4, 0.8, 1.2, 2.0m \\
Sonic anemometer                 - wind velocities (u,v,w)                 - sonic temperature & Campbell Scientific & CSAT3 & 2.5m \\
Temperature and Relative Humidity & Vaisala & HMP45A & 2.5m \\
\end{longtable}


\subsection*{Tower Global Metadata}

\begin{longtable}{p{4.5cm}p{10.5cm}}
\toprule
\textbf{Attribute} & \textbf{Value} \\
\midrule
\endhead
\bottomrule
\endfoot
\textbf{PI} & Jason Beringer \\
\textbf{Institution} & University of Western Australia \\
\textbf{Latitude} & -32.477093 \\
\textbf{Longitude} & 116.938559 \\
\textbf{Elevation} & N/A \\
\textbf{Soil Type} & ?? \\
\textbf{Vegetation} & Dry schlerophyll woodland \\
\textbf{Temporal Coverage} & 2017-10-20 13:00:00 to 2025-07-08 01:00:00 \\
\textbf{Data Link} & \url{http://data.ozflux.org.au/} \\
\end{longtable}

\newpage

\section{Site Profile: Cape Tribulation}

\subsection*{Physical Configuration \& Soil Sensor Depths}

\begin{figure}[H]
\centering
  \includegraphics[width=0.85\textwidth,height=0.45\textheight,keepaspectratio]{/home/sanjays/et97_scratch2/oldscratch/Ozflux_data_full/L6/soil_depth_plots/FluxTower_Diagrams/CapeTribulation/CapeTribulation_diagram.png}
  \caption*{Soil sensor depth profile and tower instrument heights.}
\end{figure}

\newpage

\subsection*{Tower Global Metadata}

\begin{longtable}{p{4.5cm}p{10.5cm}}
\toprule
\textbf{Attribute} & \textbf{Value} \\
\midrule
\endhead
\bottomrule
\endfoot
\textbf{PI} & Michael Liddell \\
\textbf{Institution} & James Cook University, Cairns \\
\textbf{Latitude} & -16.103219 \\
\textbf{Longitude} & 145.446922 \\
\textbf{Elevation} & N/A \\
\textbf{Soil Type} & Acidic, Dystrophic, Brown Dermosol \\
\textbf{Vegetation} & lowland complex mesophyll vine forest \\
\textbf{Temporal Coverage} & 2010-01-01 00:00:00 to 2018-11-02 12:00:00 \\
\textbf{Data Link} & \url{http://data.ozflux.org.au/} \\
\end{longtable}

\newpage

\section{Site Profile: Collie}

\subsection*{Physical Configuration \& Soil Sensor Depths}

\begin{figure}[H]
\centering
  \includegraphics[width=0.85\textwidth,height=0.45\textheight,keepaspectratio]{/home/sanjays/et97_scratch2/oldscratch/Ozflux_data_full/L6/soil_depth_plots/FluxTower_Diagrams/Collie/Collie_diagram.png}
  \caption*{Soil sensor depth profile and tower instrument heights.}
\end{figure}

\newpage

\subsection*{Site Description}

\textbf{Site Description}

The Collie flux stationis was located approximately 10km southeast of Collie, near Perth, Western Australia (GPS coordinates: -33.420, 116.237; elevation 384m)

The flux tower site was located within an area of wandoo woodland. Elevation of the site was close to 384 m. Climate information came from the nearby Collie BoM AWS station 009628 (1899 to present) and shows mean annual precipitation was approximately 933 mm with highest rainfall in June and July of \textasciitilde{}176 mm each month.

Maximum temperatures range from 30.5$^\circ$C (in Jan) to 15.5$^\circ$C (in July), while minimum temperatures range from 4.2$^\circ$C (in July) to 13.2 $^\circ$C (in Jan).

The instrument mast was 35m tall. Heat, water vapour, and carbon dioxide measurements were taken using the open-path eddy flux technique. Temperature, humidity, wind speed, wind direction, soil heat fluxes, soil temperature, and soil moisture content were also collected.

\subsection*{Purpose \& Scientific Goals}

\textbf{Purpose}

The purpose of the Collie flux station was:

\begin{itemize}
  \item to provides nationally consistent observations micrometeorology (climate, radiation, fluxes of carbon and water).
  \item to provide longterm measurements as part of the Ozflux network
\end{itemize}

The tower was a component of the UWA critical zone observatory.

\subsection*{Site Photographs}

\begin{figure}[H]
\centering
  \includegraphics[width=0.8\textwidth,height=0.4\textheight,keepaspectratio]{/home/sanjays/et97_scratch2/oldscratch/Ozflux_data_full/ozflux_mirror/www.ozflux.org.au/monitoringsites/collie/images/DJI_0025.jpg}
  \caption*{Photo 1: The Collie flux tower}
\end{figure}

\begin{figure}[H]
\centering
  \includegraphics[width=0.8\textwidth,height=0.4\textheight,keepaspectratio]{/home/sanjays/et97_scratch2/oldscratch/Ozflux_data_full/ozflux_mirror/www.ozflux.org.au/monitoringsites/collie/images/DJI_0010.jpg}
  \caption*{Photo 2: Eddy covariance instrumentation}
\end{figure}

\begin{figure}[H]
\centering
  \includegraphics[width=0.8\textwidth,height=0.4\textheight,keepaspectratio]{/home/sanjays/et97_scratch2/oldscratch/Ozflux_data_full/ozflux_mirror/www.ozflux.org.au/monitoringsites/collie/images/DJI_0002.jpg}
  \caption*{Photo 3: Meteorological instrumentation}
\end{figure}

\begin{figure}[H]
\centering
  \includegraphics[width=0.8\textwidth,height=0.4\textheight,keepaspectratio]{/home/sanjays/et97_scratch2/oldscratch/Ozflux_data_full/ozflux_mirror/www.ozflux.org.au/monitoringsites/collie/images/DJI_0028.jpg}
  \caption*{Photo 4: Looking down from above the tower}
\end{figure}

\begin{figure}[H]
\centering
  \includegraphics[width=0.8\textwidth,height=0.4\textheight,keepaspectratio]{/home/sanjays/et97_scratch2/oldscratch/Ozflux_data_full/ozflux_mirror/www.ozflux.org.au/monitoringsites/collie/images/DJI_0035.jpg}
  \caption*{Photo 5: The nearby Radio Acoustic Sounding system (RASS)}
\end{figure}

\subsection*{Measurement Instrumentation}

\begin{longtable}{p{4.5cm}p{3.5cm}p{3.5cm}p{3.5cm}}
\toprule
\textbf{Instrument Type} & \textbf{Make} & \textbf{Model} & \textbf{Situation} \\
\midrule
\endhead
\bottomrule
\endfoot
Open path  CO2, H2O         Atmospheric pressure & Li-COR & Li-7500 & 35m \\
Sonic anemometer & Campbell Scientific & CSAT3 & 35m \\
Rh+T & Vaisala & HMP45A & 35m \\
Soil Heat Flux (2  replicates) & Rebs & HFT-3 & 0.1m \\
Soil Temperature & Campbell Scientific & TCAV-L & 0.08m \\
Soil moisture & Campbell Scientific & CS650 & 0.1, 0.2, 0.6m \\
\end{longtable}


\subsection*{Tower Global Metadata}

\begin{longtable}{p{4.5cm}p{10.5cm}}
\toprule
\textbf{Attribute} & \textbf{Value} \\
\midrule
\endhead
\bottomrule
\endfoot
\textbf{PI} & Jason Beringer \\
\textbf{Institution} & University of Western Australia \\
\textbf{Latitude} & -33.42 \\
\textbf{Longitude} & 116.237 \\
\textbf{Elevation} & N/A \\
\textbf{Soil Type} & ?? \\
\textbf{Vegetation} & Dry schlerophyll woodland \\
\textbf{Temporal Coverage} & 2017-08-04 00:30:00 to 2019-11-11 13:30:00 \\
\textbf{Data Link} & \url{http://data.ozflux.org.au/} \\
\end{longtable}

\newpage

\section{Site Profile: Cow Bay}

\subsection*{Physical Configuration \& Soil Sensor Depths}

\begin{figure}[H]
\centering
  \includegraphics[width=0.85\textwidth,height=0.45\textheight,keepaspectratio]{/home/sanjays/et97_scratch2/oldscratch/Ozflux_data_full/L6/soil_depth_plots/FluxTower_Diagrams/CowBay/CowBay_diagram.png}
  \caption*{Soil sensor depth profile and tower instrument heights.}
\end{figure}

\newpage

\subsection*{Site Description}

\textbf{Site Description}

The Cow Bay flux station is located in lowland tropical rainforest at the Daintree Discovery Centre near Cow Bay in Far North Queensland (GPS coordinates: -16.238189, 145.427150; elevation: 86 m). The land is adjacent to the Daintree National Park which is part of the Wet Tropics World Heritage Area (WTWHA).

The site is flanked to the south-west by coastal ranges rising to 400 m and to the east by the Coral Sea. The red clay loam podzolic soils are of metamorphic origin and have good drainage characteristics. The Cow Bay site is gently sloping but the fetch area makes the site one of very complex terrain.

The forest is classed as complex mesophyll vine forest (type 1a) and has an average canopy height of 25 m. The dominant canopy trees belong to the Arecaceae, Euphorbiaceae, Rutaceae, Meliaceae, Myristicaceae and Icacinaceae families. It is continuous for several kilometres around the Cow Bay Tower except for an area 800 m north-east of the tower, which is cleared agricultural land used for a cattle farm. A 1 Ha census plot is located in immediate proximity to the flux tower.

Annual average rainfall at the site is around 3700 mm and is strongly seasonal, with 68\% falling between January and April (wet season).

\textbf{Daintree Discovery Centre}


          The flux station is located at the 23 m level viewing platform on the visitor tower at the Daintree Discovery Centre . The visitor centre has seen over 1 million visitors visit since it was set up in 1989 by Pam and Ron Burkett.  The DDC is recognized as a leader in ecotourism in the Daintree and has won numerous state, national and international awards including the International Award for Ecotourism (Skal Congree, Budapest 2009). In 2016 the DDC was sold to the Aboriginal Development Benefits Trust (ADBT). The ADBT helps to develop indigenous community, youth and entrepreneurship programs within the Lower Gulf region.

\subsection*{Purpose \& Scientific Goals}

\textbf{Purpose}

The purpose of the Cow Bay flux station is:

\begin{itemize}
  \item to measure exchanges of carbon dioxide, water vapour and energy 
                between the tropical lowland rainforest in Far North Queensland and 
                the atmosphere using micrometeorological techniques
  \item to quantify the changes in carbon and water balances of an Australian tropical rainforest on a long term basis in the face of climate change
  \item to present the results from the study in real time to the public and inform the public on what these results mean
\end{itemize}

This work is part of an collaborative study between JCU and the Daintree Discovery Centre.

The Cow Bay flux station is the most publicly accessible flux tower in Australia and provides considerable science outreach to the community through the Daintree Discovery Centre Visitor Information Centre.

\subsection*{Site Photographs}

\begin{figure}[H]
\centering
  \includegraphics[width=0.8\textwidth,height=0.4\textheight,keepaspectratio]{/home/sanjays/et97_scratch2/oldscratch/Ozflux_data_full/ozflux_mirror/www.ozflux.org.au/monitoringsites/cowbay/images/tower_disco_ctre1.jpg}
  \caption*{Photo 1: Cow Bay Tower (35m) EC system, view to the west}
\end{figure}

\begin{figure}[H]
\centering
  \includegraphics[width=0.8\textwidth,height=0.4\textheight,keepaspectratio]{/home/sanjays/et97_scratch2/oldscratch/Ozflux_data_full/ozflux_mirror/www.ozflux.org.au/monitoringsites/cowbay/images/tower_eddie_co.jpg}
  \caption*{Photo 2: Cow Bay Tower (23m) platform, view to the south}
\end{figure}

\subsection*{Measurement Instrumentation}

\begin{longtable}{p{4.5cm}p{3.5cm}p{3.5cm}p{3.5cm}}
\toprule
\textbf{Instrument Type} & \textbf{Make} & \textbf{Model} & \textbf{Situation} \\
\midrule
\endhead
\bottomrule
\endfoot
IRGA Open path                  - density CO2, H2O  LI-COR LI-7500RS 35m & LI-COR & LI-7500RS & 35m \\
IRGA Interface Unit                 - atmospheric pressure & LI-COR & LI-7550 & 35m \\
3D sonic anemometer                 - wind velocities (u,v,w)                  - sonic temperature & Campbell Scientific & CSAT3B & 35m \\
Air temperature, Humidity                  - T, Rh & Rotronix & HC2S3 & 35m \\
PAR Radiation                  - PPFD & LI-COR & LI-190SB & 35m \\
4 Component Radiation                  - SW-d, LW-d, LW-u, SW-u, & Hukseflux & NR01 & 35m \\
Net Radiation                  - Rn & Kipp \& Zonen & NRLite & 35m \\
Rainfall                 - Rain & Campbell Scientific & RIM8000 & 23m \\
Pyranometer & Apogee & SP-110 & 23m \\
Soil temperature                  (averaging array) & Campbell Scientific & TCAV & -7cm20m S of tower \\
Soil heat flux plates (2) & Hukseflux & HFP01 & -10cm20m S of tower \\
Soil water content                  (time domain reflectometer) & Campbell Scientific & CS616 & -6cm, -75cm, -150cm, -200cm20m S of tower \\
Data Loggers                 - EC fluxes & Campbell Scientific & CR1000X & 35m \\
Phenocam - Field Camera 5MP                 - 30 minute images & Stardot & E1159C & 35m \\
Automatic weather station                  (T, Rh, Rain) & Vaisala & WXT520 & 23m \\
\end{longtable}


\subsection*{Tower Global Metadata}

\begin{longtable}{p{4.5cm}p{10.5cm}}
\toprule
\textbf{Attribute} & \textbf{Value} \\
\midrule
\endhead
\bottomrule
\endfoot
\textbf{PI} & Michael Liddell \\
\textbf{Institution} & James Cook University, Cairns \\
\textbf{Latitude} & -16.238189 \\
\textbf{Longitude} & 145.42715 \\
\textbf{Elevation} & N/A \\
\textbf{Soil Type} & red clay loam podzolic soils \\
\textbf{Vegetation} & lowland complex mesophyll vine forest \\
\textbf{Temporal Coverage} & 2009-01-01 09:00:00 to 2025-07-23 09:00:00 \\
\textbf{Data Link} & \url{http://data.ozflux.org.au/} \\
\end{longtable}

\newpage

\section{Site Profile: Daly Pasture}

\subsection*{Physical Configuration \& Soil Sensor Depths}

\begin{figure}[H]
\centering
  \includegraphics[width=0.85\textwidth,height=0.45\textheight,keepaspectratio]{/home/sanjays/et97_scratch2/oldscratch/Ozflux_data_full/L6/soil_depth_plots/FluxTower_Diagrams/DalyPasture/DalyPasture_diagram.png}
  \caption*{Soil sensor depth profile and tower instrument heights.}
\end{figure}

\newpage

\subsection*{Site Description}

\textbf{Site Description}

The Daly River Pasture flux station was located 62km south west of Pine Creek, Northern Territory (GPS coordinates: -14.0633, 131.3181).

The flux  tower site was identified as tropical pasture, and the vegetation was dominated by species \textit{Chamaecrista  rotundifolia}(Round-leaf cassia cv. Wynn), \textit{Digitaria milijiana} (Jarra grass) and \textit{Aristida sp}. standing at approximately 0.3m tall. The soil at the site was a mixture of red kandosol and deep sand. Elevation of the site was close to 70m and mean annual precipitation at a nearby Bureau of Meteorology site is 1250mm.

Maximum temperatures ranged from 37.5$^\circ$C (in October) to 31.2$^\circ$C  (in June), while minimum temperatures ranged from 12.6$^\circ$C (in July) to  23.8$^\circ$C (in January). Maximum temperatures varied on a seasonal basis between 6.3$^\circ$C while minimum temperatures varied by 11.2$^\circ$C.

The instrument mast was 15 meters tall. Heat, water vapour and carbon dioxide measurements were taken using the open-path eddy flux technique. Temperature, humidity, wind speed, wind  direction, rainfall, incoming and reflected shortwave radiation and net  radiation were measured.

Ancillary measurements taken at the site included LAI, leaf-scale physiological properties (gas exchange, leaf isotope ratios, N and chlorophyll concentrations), vegetation optical  properties and soil physical properties. Airborne based remote sensing (Lidar  and hyperspectral measurements) was carried out across the transect in  September 2008.

\subsection*{Purpose \& Scientific Goals}

\textbf{Purpose}

The  purpose of the Daly River Pasture Flux station was to:

\begin{itemize}
  \item Provide  information as part of a larger network of flux towers established along the  North Australian Tropical Transect (NATT) gradient, which extends \textasciitilde{}1000km  south from Darwin 12.5$^\circ$S.  The towers provide flux data across the savanna’s heterogeneous ‘top end’  ecosystems including open forest savanna, open savanna woodland and seasonally  inundated floodplains.
  \item Examine spatial patterns and processes of  land-surface-atmosphere exchanges (radiation, heat, moisture, CO2 and other  trace gasses) across scales from leaf to landscape scales within Australian  savannas.
  \item Determine  the climate and ecosystem characteristics (physical structure,  species composition, physiological function) that drive spatial and temporal variations of  carbon, water and energy fluxes from north Australian savanna.
  \item Determine if fluxes of carbon, water vapour and  heat over the various ecosystems as derived from the various measurement  techniques can be combined to form a comprehensive and consistent estimate of  the regional fluxes and budgets across the landscape.
  \item Quantify  fluxes over land use which includes grazing, introduced pastures and low stock  density.
  \item Provide  longer term measurements for future projects.
\end{itemize}

\subsection*{Measurement Instrumentation}

\begin{longtable}{p{4.5cm}p{3.5cm}p{3.5cm}p{3.5cm}}
\toprule
\textbf{Instrument Type} & \textbf{Make} & \textbf{Model} & \textbf{Location} \\
\midrule
\endhead
\bottomrule
\endfoot
Albedometer & Kipp and Zonen & CM7B & 15m \\
Data  logger & Campbell Scientific & CR3000 & Ground \\
Infra-red  gas analyser & Licor & LI7500 & 15m \\
Net  radiometer & Kipp and Zonen & NRlite & 15m \\
Rain gauge & Hydrological services & CS702 & 15m \\
Soil  heat flux (two replicates) & Campbell Scientific & HFT3 & Ground,  -0.15m \\
Soil  moisture (two replicates) & Campbell Scientific & CS616 & -0.05,  0.50m \\
Solar  regulator & ProStar & PS-30M & Ground \\
Sonic anemometer       - wind velocities (u,v,w)        - sonic temperature & Campbell Scientific & CSAT3 & 15m \\
Temperature  and Relative Humidity & Vaisala & HMP-45AC & 15m \\
Unidirectional  pyrgeometer & Kipp and Zonen & CG-2 & 15m \\
\end{longtable}


\subsection*{Tower Global Metadata}

\begin{longtable}{p{4.5cm}p{10.5cm}}
\toprule
\textbf{Attribute} & \textbf{Value} \\
\midrule
\endhead
\bottomrule
\endfoot
\textbf{PI} & Jason Beringer and Lindsay Hutley \\
\textbf{Institution} & University of Western Australia and Charles Darwin University \\
\textbf{Latitude} & -14.0633 \\
\textbf{Longitude} & 131.3181 \\
\textbf{Elevation} & N/A \\
\textbf{Soil Type} & Red Kandosol \\
\textbf{Vegetation} & Tropical improved pasture \\
\textbf{Temporal Coverage} & 2008-01-01 00:00:00 to 2013-09-08 13:30:00 \\
\textbf{Data Link} & \url{http://data.ozflux.org.au/} \\
\end{longtable}

\newpage

\section{Site Profile: Daly Uncleared}

\subsection*{Physical Configuration \& Soil Sensor Depths}

\begin{figure}[H]
\centering
  \includegraphics[width=0.85\textwidth,height=0.45\textheight,keepaspectratio]{/home/sanjays/et97_scratch2/oldscratch/Ozflux_data_full/L6/soil_depth_plots/FluxTower_Diagrams/DalyUncleared/DalyUncleared_diagram.png}
  \caption*{Soil sensor depth profile and tower instrument heights.}
\end{figure}

\newpage

\subsection*{Site Description}

\textbf{Site Description}

The Daly  River Uncleared flux tower site is located in the Douglas River Daly River  Esplanade Conservation area, approximately 60 km south west of Pine Creek,  Northern Territory (GPS coordinates: -14.1592, 131.3881).

The flux tower site is classified as a Woodland savanna. The overstory is co dominated by tree species \textit{E. tetrodonta, C. latifolia,} \textit{Terminalia grandiflora, Sorghum  sp. and Heteropogon triticeus.} Average canopy height measures 16.4 m. Elevation of the site is close to 110m and mean annual  precipitation at a nearby Bureau of Meteorology site is 1170mm.

Maximum  temperatures range from 37.5$^\circ$C (in October) to 31.2$^\circ$C (in  June), while minimum temperatures range from 12.6$^\circ$C (in July) to  23.8$^\circ$C (in January). Maximum temperatures range seasonally by 6.3$^\circ$C and minimum temperatures by 11.2$^\circ$C.

The instrument mast is 23 meters  tall. Heat, water vapour and carbon dioxide measurements are taken using the  open-path eddy flux technique. Temperature, humidity, wind speed, wind  direction, rainfall, incoming and reflected shortwave radiation and net  radiation are measured above the canopy.

Ancillary measurements taken at  the site include LAI, leaf-scale physiological properties (gas exchange, leaf  isotope ratios, N and chlorophyll concentrations), vegetation optical  properties and soil physical properties. Airborne based remote sensing (Lidar  and hyperspectral measurements) was carried out across the transect in  September 2008.

\subsection*{Purpose \& Scientific Goals}

\textbf{Purpose}

The  primary purpose of the Daly River Uncleared Flux Station is to:

\begin{itemize}
  \item Provide  information as part of a larger network of flux towers established along the  North Australian Tropical Transect (NATT) gradient, which extends \textasciitilde{}1000km  south from Darwin 12.5$^\circ$S.  The towers provide flux data across the savanna’s heterogeneous ‘top end’  ecosystems including open forest savanna, open savanna woodland and seasonally  inundated floodplains.
  \item Examine spatial patterns and processes of  land-surface-atmosphere exchanges (radiation, heat, moisture, CO2 and other  trace gasses) across scales from leaf to landscape scales within Australian  savannas.
  \item Determine  the climate and ecosystem characteristics (physical structure,  species composition, physiological function) that drive spatial and temporal variations of  carbon, water and energy fluxes from north Australian savanna.
  \item Determine if fluxes of carbon, water vapour and  heat over the various ecosystems as derived from the various measurement  techniques can be combined to form a comprehensive and consistent estimate of  the regional fluxes and budgets across the landscape.
  \item Quantify  fluxes over woodland savanna.
  \item Provide  longer term measurements for future projects.
\end{itemize}

\subsection*{Site Photographs}

\begin{figure}[H]
\centering
  \includegraphics[width=0.8\textwidth,height=0.4\textheight,keepaspectratio]{/home/sanjays/et97_scratch2/oldscratch/Ozflux_data_full/ozflux_mirror/www.ozflux.org.au/monitoringsites/dalyuncleared/images/DalyUC1.JPG}
  \caption*{Photo 1: View  of canopy from top of instrument mast}
\end{figure}

\begin{figure}[H]
\centering
  \includegraphics[width=0.8\textwidth,height=0.4\textheight,keepaspectratio]{/home/sanjays/et97_scratch2/oldscratch/Ozflux_data_full/ozflux_mirror/www.ozflux.org.au/monitoringsites/dalyuncleared/images/DalyUC2.jpg}
  \caption*{Photo 2: Taking  ancillary measurements during the SPECIAL campaign, September 2008}
\end{figure}

\begin{figure}[H]
\centering
  \includegraphics[width=0.8\textwidth,height=0.4\textheight,keepaspectratio]{/home/sanjays/et97_scratch2/oldscratch/Ozflux_data_full/ozflux_mirror/www.ozflux.org.au/monitoringsites/dalyuncleared/images/DalyUC3.jpg}
  \caption*{Photo 3: View of  23m tall mast above the woody savanna}
\end{figure}

\begin{figure}[H]
\centering
  \includegraphics[width=0.8\textwidth,height=0.4\textheight,keepaspectratio]{/home/sanjays/et97_scratch2/oldscratch/Ozflux_data_full/ozflux_mirror/www.ozflux.org.au/monitoringsites/dalyuncleared/images/DalyUC4.JPG}
  \caption*{Photo 4: The  understory of the tower site}
\end{figure}

\subsection*{Measurement Instrumentation}

\begin{longtable}{p{4.5cm}p{3.5cm}p{3.5cm}p{3.5cm}}
\toprule
\textbf{Instrument Type} & \textbf{Make} & \textbf{Model} & \textbf{Location} \\
\midrule
\endhead
\bottomrule
\endfoot
Bidirectional  pyranometer / pyrgeometer & Kipp and Zonen & CNR4 & 21m \\
Data  logger & Campbell Scientific & CR3000 & Ground \\
Infra-red  gas analyser & Licor & LI7500 & 23m \\
Net  radiometer & Kipp and Zonen & NRlite & 21m \\
Rain gauge & Hydrological services & CS702 & Ground \\
Soil  heat flux & Campbell Scientific & HFT3 & -0.15m \\
Soil  moisture (two replicates) & Campbell Scientific & CS616 & -0.05,  0.50m \\
Solar  regulator & ProStar & PS-30M & Ground \\
Sonic anemometer       - wind velocities (u,v,w)        - sonic temperature & Campbell Scientific & CSAT3 & 23m \\
Temperature  and Relative Humidity  (two replicates) & Vaisala & HMP-45AC & 2m, 23m \\
\end{longtable}


\subsection*{Tower Global Metadata}

\begin{longtable}{p{4.5cm}p{10.5cm}}
\toprule
\textbf{Attribute} & \textbf{Value} \\
\midrule
\endhead
\bottomrule
\endfoot
\textbf{PI} & Lindsay Hutley and Jason Beringer \\
\textbf{Institution} & University of Western Australia \\
\textbf{Latitude} & -14.1592 \\
\textbf{Longitude} & 131.3881 \\
\textbf{Elevation} & N/A \\
\textbf{Soil Type} & red kandasol \\
\textbf{Vegetation} & Woody Savanna \\
\textbf{Temporal Coverage} & 2008-01-01 00:00:00 to 2025-07-07 15:30:00 \\
\textbf{Data Link} & \url{http://data.ozflux.org.au/} \\
\end{longtable}

\newpage

\section{Site Profile: DigbyPlantation}

\subsection*{Physical Configuration \& Soil Sensor Depths}

\begin{figure}[H]
\centering
  \includegraphics[width=0.85\textwidth,height=0.45\textheight,keepaspectratio]{/home/sanjays/et97_scratch2/oldscratch/Ozflux_data_full/L6/soil_depth_plots/FluxTower_Diagrams/DigbyPlantation/DigbyPlantation_diagram.png}
  \caption*{Soil sensor depth profile and tower instrument heights.}
\end{figure}

\newpage

\subsection*{Tower Global Metadata}

\begin{longtable}{p{4.5cm}p{10.5cm}}
\toprule
\textbf{Attribute} & \textbf{Value} \\
\midrule
\endhead
\bottomrule
\endfoot
\textbf{PI} & Edoardo Daly \\
\textbf{Institution} &  \\
\textbf{Latitude} & -37.882694 \\
\textbf{Longitude} & 141.616472 \\
\textbf{Elevation} & N/A \\
\textbf{Soil Type} &  \\
\textbf{Vegetation} & Eucalyptus globulus (blue gum) \\
\textbf{Temporal Coverage} & 2018-01-01 00:30:00 to 2021-07-07 12:00:00 \\
\textbf{Data Link} & \url{http://data.ozflux.org.au/} \\
\end{longtable}

\newpage

\section{Site Profile: Fogg Dam}

\subsection*{Physical Configuration \& Soil Sensor Depths}

\begin{figure}[H]
\centering
  \includegraphics[width=0.85\textwidth,height=0.45\textheight,keepaspectratio]{/home/sanjays/et97_scratch2/oldscratch/Ozflux_data_full/L6/soil_depth_plots/FluxTower_Diagrams/FoggDam/FoggDam_diagram.png}
  \caption*{Soil sensor depth profile and tower instrument heights.}
\end{figure}

\newpage

\subsection*{Site Description}

\textbf{Site Description}

The Fogg  Dam flux station was located approximately 6km east of Black Jungle, Northern  Territory (GPS coordinates: -12.5452, 131.3072).

The flux tower site was classified as a seasonally  flooded wetland. The vegetation was dominated by species \textit{Oryza rufipogon, Pseudoraphis spinescens and Eleocharis dulcis}. Elevation of the site was close to 4m and mean annual  precipitation from a nearby Bureau of Meteorology site measured 1411mm.

Maximum temperatures ranged from 31.3$^\circ$C (in June and July) to 35.6$^\circ$C  (in October), while minimum temperatures ranged from 14.9$^\circ$C (in  July) to 23.9$^\circ$C (in December and February). Maximum temperatures  varied on a seasonal basis by approximately 4.3$^\circ$C and minimum  temperatures by 9.0$^\circ$C.

The instrument mast was 15m tall.  Heat, water vapour and carbon dioxide measurements are taken using the  open-path eddy flux technique. Temperature, humidity, wind speed, wind  direction, rainfall, incoming and reflected shortwave radiation and net  radiation were measured above the canopy. Soil heat fluxes were measured and  soil moisture content was gathered using time domain reflectometry.

Ancillary measurements taken at  the site include LAI, leaf-scale physiological properties (gas exchange, leaf  isotope ratios, N and chlorophyll concentrations), vegetation optical  properties and soil physical properties. Airborne based remote sensing (Lidar  and hyperspectral measurements) was carried out across the transect in  September 2008.

\subsection*{Purpose \& Scientific Goals}

\textbf{Purpose}

The primary purpose of the Fogg Dam flux station was to:

\begin{itemize}
  \item Provide information as part of a larger network of flux towers established along the North Australian Tropical Transect (NATT) gradient, which extends \textasciitilde{}1000km south from Darwin 12.5$^\circ$ S. The towers provide flux data across the savanna's heterogeneous 'top end' ecosystems including open forest savanna, open savanna woodland and seasonally inundated floodplains.
  \item Examine spatial patterns and processes of land-surface-atmosphere exchanges (radiation, heat, moisture, CO2 and other trace gasses) across scales from leaf to landscape scales within Australian savannas.
  \item Determine the climate and ecosystem characteristics (physical structure, species composition, physiological function) that drive spatial and temporal variations of carbon, water and energy fluxes from north Australian savanna.
  \item Determine if fluxes of carbon, water vapour and heat over the various ecosystems as derived from the various measurement techniques can be combined to form a comprehensive and consistent estimate of the regional fluxes and budgets across the landscape.
\end{itemize}

\subsection*{Site Photographs}

\begin{figure}[H]
\centering
  \includegraphics[width=0.8\textwidth,height=0.4\textheight,keepaspectratio]{/home/sanjays/et97_scratch2/oldscratch/Ozflux_data_full/ozflux_mirror/www.ozflux.org.au/monitoringsites/foggdam/images/FD1.JPG}
  \caption*{Photo 1: Fogg Dam   Flux Tower}
\end{figure}

\begin{figure}[H]
\centering
  \includegraphics[width=0.8\textwidth,height=0.4\textheight,keepaspectratio]{/home/sanjays/et97_scratch2/oldscratch/Ozflux_data_full/ozflux_mirror/www.ozflux.org.au/monitoringsites/foggdam/images/FD2.JPG}
  \caption*{Photo 2: LiCor-7500 and CSAT3 mounted on the tower}
\end{figure}

\begin{figure}[H]
\centering
  \includegraphics[width=0.8\textwidth,height=0.4\textheight,keepaspectratio]{/home/sanjays/et97_scratch2/oldscratch/Ozflux_data_full/ozflux_mirror/www.ozflux.org.au/monitoringsites/foggdam/images/FD3.JPG}
  \caption*{Photo 3: Net and Four component Kipp and Zonen radiometers mounted on the  tower}
\end{figure}

\subsection*{Measurement Instrumentation}

\begin{longtable}{p{4.5cm}p{3.5cm}p{3.5cm}p{3.5cm}}
\toprule
\textbf{Instrument Type} & \textbf{Make} & \textbf{Model} & \textbf{Situation} \\
\midrule
\endhead
\bottomrule
\endfoot
Pyrgeometer & Kipp \& Zonen & CG4 & 14.5m \\
Data Logger & Campbell Scientific & CR3000 & Ground \\
Infra-red  gas analyser & Licor & LI7500 & 14.5m \\
Net  radiometer & Kipp \& Zonen & N-R  Lite & Ground \\
Rain  gauge & Campbell Scientific & CS702 & 45m \\
Soil Heat Flux               (two replicates) & Campbell Scientific & HFT3 & Ground,  -0.15m \\
Soil  Moisture & Campbell Scientific & CS616 & -0.50m \\
Sonic anemometer                 - wind velocities (u,v,w)                - sonic temperature & Campbell Scientific & CSAT3 & 14.5m \\
Temperature and Relative Humidity probe & Vaisalla & HMP45A & 14.5m \\
\end{longtable}


\subsection*{Tower Global Metadata}

\begin{longtable}{p{4.5cm}p{10.5cm}}
\toprule
\textbf{Attribute} & \textbf{Value} \\
\midrule
\endhead
\bottomrule
\endfoot
\textbf{PI} & Jason Beringer and Lindsay Hutley \\
\textbf{Institution} & Monash University and Charles Darwin University \\
\textbf{Latitude} & -12.5452 \\
\textbf{Longitude} & 131.3072 \\
\textbf{Elevation} & N/A \\
\textbf{Soil Type} & rich organic \\
\textbf{Vegetation} & Ephemeral tropical wetland \\
\textbf{Temporal Coverage} & 2006-02-07 00:00:00 to 2008-10-31 17:30:00 \\
\textbf{Data Link} & \url{http://data.ozflux.org.au/} \\
\end{longtable}

\newpage

\section{Site Profile: GatumPasture}

\subsection*{Physical Configuration \& Soil Sensor Depths}

\begin{figure}[H]
\centering
  \includegraphics[width=0.85\textwidth,height=0.45\textheight,keepaspectratio]{/home/sanjays/et97_scratch2/oldscratch/Ozflux_data_full/L6/soil_depth_plots/FluxTower_Diagrams/GatumPasture/GatumPasture_diagram.png}
  \caption*{Soil sensor depth profile and tower instrument heights.}
\end{figure}

\newpage

\subsection*{Site Description}

\textbf{Site Description}

The Gatum Pasture pasture flux station was located on private farmland in south western Victoria, approximately 260km west of Melbourne and 15km north west of the nearest settlement, Cavendish (coordinates: -37.3903, 141.9716). Elevation of the site is 250m asl, and the local terrain has a very slight (<1$^\circ$) south easterly aspect.

The ecosystem was a sown pasture dominated by phalaris (\textit{Phalaris aquatica L.}) and subterrenean clover (\textit{Trifolium subteraneum L.}) that was regularly fertilised with superphosphate. The pasture was intensively grazed by sheep and cattle. The grass canopy was thus generally <10cm tall, and largely died off during summer and early autumn. Mean annual precipitation was 736mm (Cavendish Bureau of Meteorology station [ID 089043], 1953-2021), and mean annual temperature was 19.2$^\circ$C (Hamilton Airport Bureau of Meteorology station [ID 090173], 1983-2021).

The instrument mast was 3.5m tall. Fluxes of heat, water vapour and carbon dioxide were measured using the open-path eddy covariance technique (at height of 3m). Supplementary measurements above the canopy included temperature, humidity, windspeed, wind direction, rainfall, incoming and reflected shortwave radiation and net radiation. Soil moisture content was measured using frequency domain reflectometry, while soil heat fluxes and temperature were also measured.

\subsection*{Purpose \& Scientific Goals}

\textbf{Purpose}

The purpose of the Gatum Pasture flux station was to:

\begin{itemize}
  \item identify differences in water use between pastures and plantations.
  \item quantify evapotranspiration for use in and validation of simple catchment-scale hydrological models.
  \item quantify long-term carbon exchanges (photosynthesis and respiration) between ecosystem and atmosphere.
\end{itemize}

\subsection*{Site Photographs}

\begin{figure}[H]
\centering
  \includegraphics[width=0.8\textwidth,height=0.4\textheight,keepaspectratio]{/home/sanjays/et97_scratch2/oldscratch/Ozflux_data_full/ozflux_mirror/www.ozflux.org.au/monitoringsites/gatumpasture/images/IMG_2796.jpg}
  \caption*{Photo 1: Flux tower and instrumentation}
\end{figure}

\begin{figure}[H]
\centering
  \includegraphics[width=0.8\textwidth,height=0.4\textheight,keepaspectratio]{/home/sanjays/et97_scratch2/oldscratch/Ozflux_data_full/ozflux_mirror/www.ozflux.org.au/monitoringsites/gatumpasture/images/IMG_2797.jpg}
  \caption*{Photo 2: Surrounding pasture}
\end{figure}

\subsection*{Measurement Instrumentation}

\begin{longtable}{p{4.5cm}p{3.5cm}p{3.5cm}p{3.5cm}}
\toprule
\textbf{Instrument Type} & \textbf{Make} & \textbf{Model} & \textbf{Situation} \\
\midrule
\endhead
\bottomrule
\endfoot
Open path  CO2, H2O & Campbell Scientific & IRGASON & Mount arm (W), 3m \\
Sonic anemometer & Campbell Scientific & IRGASON & Mount arm (W), 3m \\
Net Radiometer & Kipp and Zonen & CNR4 & Mount arm (E), 3m \\
Rh+T & Vaisala & HMP110 & Mast, 3m \\
Rainfall & Hydrological Services & TB3 & Separate post \\
Soil Temperature & Campbell Scientific & TCAV & Ground, -2cm, -6cm \\
Soil Heat Flux & Huxeflux & HFP01 & Ground, -8cm \\
Soil moisture & Campbell Scientific & CS650 & Ground, -5 (3 reps), -15 (3 reps), -30 (2 reps), -50, -70, -90, -130, -150cm \\
\end{longtable}


\subsection*{Tower Global Metadata}

\begin{longtable}{p{4.5cm}p{10.5cm}}
\toprule
\textbf{Attribute} & \textbf{Value} \\
\midrule
\endhead
\bottomrule
\endfoot
\textbf{PI} & Edoardo Daly \\
\textbf{Institution} & Monash University \\
\textbf{Latitude} & -37.390033 \\
\textbf{Longitude} & 141.960897 \\
\textbf{Elevation} & N/A \\
\textbf{Soil Type} &  \\
\textbf{Vegetation} & Phalaris and Subterranean Clover \\
\textbf{Temporal Coverage} & 2015-02-09 00:30:00 to 2021-12-31 23:30:00 \\
\textbf{Data Link} & \url{http://data.ozflux.org.au/} \\
\end{longtable}

\newpage

\section{Site Profile: Gingin}

\subsection*{Physical Configuration \& Soil Sensor Depths}

\begin{figure}[H]
\centering
  \includegraphics[width=0.85\textwidth,height=0.45\textheight,keepaspectratio]{/home/sanjays/et97_scratch2/oldscratch/Ozflux_data_full/L6/soil_depth_plots/FluxTower_Diagrams/Gingin/Gingin_diagram.png}
  \caption*{Soil sensor depth profile and tower instrument heights.}
\end{figure}

\newpage

\subsection*{Site Description}

\textbf{Site Description}

The Gnangara flux station is located in coastal heath \textit{Banksia} woodland on the Swan Coastal Plain 70km north of Perth, Western Australia: 
(GPS coordinates: -31.3764, 115.7139; elevation: 51m),   and 2km south of the University   of Western Australia  International  Gravity Wave Observatory .

Climate summary taken from SILO DataDrill .

The site is a natural woodland of high species diversity. The overstorey is dominated by \textit{Banksia spp}. mainly \textit{B. menziesii, B. attenuata,}and\textit{B. grandis} with a height of around 7m and leaf area index of about 0.8.

There are occasional stands of eucalypts and  acacia that reach to 10m and have a denser foliage cover.

There are many former wetlands dotted around the  woodland, most of which were inundated all winter and some had permanent water  30 years ago. The watertable has now fallen below the base of these systems and  they are disconnected and are no longer permanently wet. The fine sediments,  sometimes diatomaceous, hold water and they have perched watertables each  winter.  There is a natural progression  of species accompanying this process as they gradually become more dominated by  more xeric species.

The soils are mainly Podosol sands, with low moisture  holding capacity. Field capacity typically about 8 to 10\%, and in summer these  generally hold less than 2\% moisture.

The watertable is at about 8.5 m below the surface, and a WA Dept of water long-term monitoring piezometer is near the base of the tower.

The instrument mast is 14m tall, with the eddy covariance instruments mounted at 14.8m.

Fluxes of carbon dioxide, water vapour and heat  are quantified with open-path eddy covariance instrumentation.

Ancillary measurements include temperature, air  humidity, wind speed and direction, precipitation, incoming and outgoing  shortwave radiation, incoming and outgoing long wave radiation, incoming total  and diffuse PAR and reflected PAR.

Soil water content and temperature are measured  at six soil depths.  Surface soil heat fluxes are also measured.

A COSMOS Cosmic ray soil moisture instrument is  installed, along with a logged piezometer, and nested piezometers installed  with short screens for groundwater profile sampling.

To monitor the watertable gradient, piezometers will be installed 500 m esat and west of the tower.

\subsection*{Purpose \& Scientific Goals}

\textbf{Purpose}

The purpose of the Gingin   flux station is:

\begin{itemize}
  \item To quantify recharge to Gnangara groundwater  mound, Perth’s most important water resource
  \item Monitor ecophysiological responses to long-term  variation in climate and water table drawdown
  \item To quantify landscape-scale exchange of carbon  dioxide, water vapour and energy in a coastal heath environment
  \item To further understand groundwater recharge under  changing climate
  \item To provide ecophysiological and  micrometeorological data representative of an important biome within Australia  subject to drying climate, falling watertables, fire and encroachment of feral  species
  \item To provide enhanced datasets of landscape-scale  exchange of carbon dioxide, water vapour and energy along with ecophysiological  characteristics and drivers in a semi-arid temperate ecosystems in Australia.
\end{itemize}

\subsection*{Site Photographs}

\begin{figure}[H]
\centering
  \includegraphics[width=0.8\textwidth,height=0.4\textheight,keepaspectratio]{/home/sanjays/et97_scratch2/oldscratch/Ozflux_data_full/ozflux_mirror/www.ozflux.org.au/monitoringsites/gingin/images/GnangaraBanksia.jpg}
  \caption*{Photo 1: Banksia woodland at the site of the tower}
\end{figure}

\begin{figure}[H]
\centering
  \includegraphics[width=0.8\textwidth,height=0.4\textheight,keepaspectratio]{/home/sanjays/et97_scratch2/oldscratch/Ozflux_data_full/ozflux_mirror/www.ozflux.org.au/monitoringsites/gingin/images/human-protractor.jpg}
  \caption*{Photo 2: The tower goes up at Gingin: June, 2011  Grant the human protractor, measuring exactly 120$^\circ$ between guy wires.}
\end{figure}

\begin{figure}[H]
\centering
  \includegraphics[width=0.8\textwidth,height=0.4\textheight,keepaspectratio]{/home/sanjays/et97_scratch2/oldscratch/Ozflux_data_full/ozflux_mirror/www.ozflux.org.au/monitoringsites/gingin/images/hole-shape-wrong.jpg}
  \caption*{Photo 3: …  “The hole’s the wrong bloody shape”}
\end{figure}

\begin{figure}[H]
\centering
  \includegraphics[width=0.8\textwidth,height=0.4\textheight,keepaspectratio]{/home/sanjays/et97_scratch2/oldscratch/Ozflux_data_full/ozflux_mirror/www.ozflux.org.au/monitoringsites/gingin/images/CoffeeRock1.jpg}
  \caption*{Photo 4: We struck "coffee rock" – iron cemented sand – at 1.20m in one of the three guy wire holes. Hard on top, softer underneath, about 50cm thick, roots do penetrate.}
\end{figure}

\begin{figure}[H]
\centering
  \includegraphics[width=0.8\textwidth,height=0.4\textheight,keepaspectratio]{/home/sanjays/et97_scratch2/oldscratch/Ozflux_data_full/ozflux_mirror/www.ozflux.org.au/monitoringsites/gingin/images/CoffeeRock2.jpg}
  \caption*{Photo 5: We struck "coffee rock" – iron cemented sand – at 1.20m in one of the three guy wire holes. Hard on top, softer underneath, about 50cm thick, roots do penetrate.}
\end{figure}

\begin{figure}[H]
\centering
  \includegraphics[width=0.8\textwidth,height=0.4\textheight,keepaspectratio]{/home/sanjays/et97_scratch2/oldscratch/Ozflux_data_full/ozflux_mirror/www.ozflux.org.au/monitoringsites/gingin/images/guy-wire-footing.jpg}
  \caption*{Photo 6: Guy wire anchor.}
\end{figure}

\begin{figure}[H]
\centering
  \includegraphics[width=0.8\textwidth,height=0.4\textheight,keepaspectratio]{/home/sanjays/et97_scratch2/oldscratch/Ozflux_data_full/ozflux_mirror/www.ozflux.org.au/monitoringsites/gingin/images/in-charge.jpg}
  \caption*{Photo 7: Somebody has to be in  charge.}
\end{figure}

\begin{figure}[H]
\centering
  \includegraphics[width=0.8\textwidth,height=0.4\textheight,keepaspectratio]{/home/sanjays/et97_scratch2/oldscratch/Ozflux_data_full/ozflux_mirror/www.ozflux.org.au/monitoringsites/gingin/images/empty-chair.jpg}
  \caption*{Photo 8: We can deal with them.}
\end{figure}

\begin{figure}[H]
\centering
  \includegraphics[width=0.8\textwidth,height=0.4\textheight,keepaspectratio]{/home/sanjays/et97_scratch2/oldscratch/Ozflux_data_full/ozflux_mirror/www.ozflux.org.au/monitoringsites/gingin/images/hole.jpg}
  \caption*{Photo 9: We can deal with them.}
\end{figure}

\begin{figure}[H]
\centering
  \includegraphics[width=0.8\textwidth,height=0.4\textheight,keepaspectratio]{/home/sanjays/et97_scratch2/oldscratch/Ozflux_data_full/ozflux_mirror/www.ozflux.org.au/monitoringsites/gingin/images/Tower-Install-022.jpg}
  \caption*{Photo 10: Plus or minus 2$^\circ$  from vertical}
\end{figure}

\begin{figure}[H]
\centering
  \includegraphics[width=0.8\textwidth,height=0.4\textheight,keepaspectratio]{/home/sanjays/et97_scratch2/oldscratch/Ozflux_data_full/ozflux_mirror/www.ozflux.org.au/monitoringsites/gingin/images/Tower-Install-026.jpg}
  \caption*{Photo 11: The first section}
\end{figure}

\begin{figure}[H]
\centering
  \includegraphics[width=0.8\textwidth,height=0.4\textheight,keepaspectratio]{/home/sanjays/et97_scratch2/oldscratch/Ozflux_data_full/ozflux_mirror/www.ozflux.org.au/monitoringsites/gingin/images/Tower-Install-043.jpg}
  \caption*{Photo 12: The second section}
\end{figure}

\begin{figure}[H]
\centering
  \includegraphics[width=0.8\textwidth,height=0.4\textheight,keepaspectratio]{/home/sanjays/et97_scratch2/oldscratch/Ozflux_data_full/ozflux_mirror/www.ozflux.org.au/monitoringsites/gingin/images/Tower-Install-045.jpg}
  \caption*{Photo 13: The last section}
\end{figure}

\begin{figure}[H]
\centering
  \includegraphics[width=0.8\textwidth,height=0.4\textheight,keepaspectratio]{/home/sanjays/et97_scratch2/oldscratch/Ozflux_data_full/ozflux_mirror/www.ozflux.org.au/monitoringsites/gingin/images/Tower-Install-053.jpg}
  \caption*{Photo 14: And she’s up}
\end{figure}

\subsection*{Measurement Instrumentation}

\begin{longtable}{p{4.5cm}p{3.5cm}p{3.5cm}p{3.5cm}}
\toprule
\textbf{Instrument Type} & \textbf{Make} & \textbf{Model} & \textbf{Location} \\
\midrule
\endhead
\bottomrule
\endfoot
Open path         CO2, H2O & LI-Cor & LI-7500 & 15m \\
3D Sonic anemometer & Campbell Scientific & CSAT3 & 15m \\
Air temperature & Vaisala & HMP45C & 5m, 15m \\
Air relative humidity & Vaisala & HMP45C & 5m, 15m \\
Net radiometer & Kipp and Zonen & CNR1 & 15m \\
Wind  speed and direction & Gill & Windsonic4-L 2D sonic & 15m \\
Precipitation & Hydrological services & CS701 & 0m x 3 \\
Soil Heat Flux plate (2  replicates) & Middleton & CN3 & -0.02m each of 3 \\
Soil temperature averaging probe       (2 replicates) & Campbell Scientific & TCAV & -0.02 $\pm$ 0.05m each of 3 \\
Soil temperature thermistor & Campbell Scientific & CS107 & -0.10, -0.20, -0.40, -0.80,  -1.80m \\
Soil water content  reflectometers (4x4depths) & Campbell Scientific & CS616 & -0.10 (x4),  -0.20 (x2),  -0.40 (x4),  -0.80 (x2),  -1.6 (x4)m \\
Soil  moisture (method 1) &  & COSMOS & 100m south-east of tower \\
Soil  moisture (method 2) & Campbell & Neutron probe & Every 25cm from the surface to 12m,       3 holes \\
Watertable & Piezometer installation,  In situ pressure transducer &  & -15m,       3 piezometers, one at site, one 400m east and west \\
\end{longtable}


\subsection*{Tower Global Metadata}

\begin{longtable}{p{4.5cm}p{10.5cm}}
\toprule
\textbf{Attribute} & \textbf{Value} \\
\midrule
\endhead
\bottomrule
\endfoot
\textbf{PI} & Richard Silberstein \\
\textbf{Institution} & Edith Cowan University, Centre for Ecosystem Management \\
\textbf{Latitude} & -31.376351 \\
\textbf{Longitude} & 115.71377 \\
\textbf{Elevation} & N/A \\
\textbf{Soil Type} & Podosol, Gutless sand \\
\textbf{Vegetation} & Banksia woodland \\
\textbf{Temporal Coverage} & 2011-10-13 14:00:00 to 2025-06-27 23:30:00 \\
\textbf{Data Link} & \url{http://data.ozflux.org.au/} \\
\end{longtable}

\newpage

\section{Site Profile: Howard Springs}

\subsection*{Physical Configuration \& Soil Sensor Depths}

\begin{figure}[H]
\centering
  \includegraphics[width=0.85\textwidth,height=0.45\textheight,keepaspectratio]{/home/sanjays/et97_scratch2/oldscratch/Ozflux_data_full/L6/soil_depth_plots/FluxTower_Diagrams/HowardSprings/HowardSprings_diagram.png}
  \caption*{Soil sensor depth profile and tower instrument heights.}
\end{figure}

\newpage

\subsection*{Site Description}

\textbf{Site Description}

The Howard Springs flux station is located in the Black Jungle Conservation Reserve in the Northern Territory (GPS coordinates: -12.4943, 131.1523).

The flux tower site is classified as an open woodland savanna. The overstory is co-dominated by tree species \textit{Eucalyptus miniata} and \textit{Eucalyptus tentrodonata}, and average tree height is 14–16m.

Elevation of the site is close to 64m and mean annual precipitation is 1750mm. Maximum temperatures range from 30.4$^\circ$C (in July) to 33.2$^\circ$C (in November), while minimum temperatures range from 19.3$^\circ$C (in July) to 25.4$^\circ$C (in November). Therefore, the maximum and minimum range varies from 7$^\circ$C (wet season) to 11$^\circ$C (dry season).

The instrument mast is 23m tall. Heat, water vapour and carbon dioxide measurements are taken using the open-path eddy flux technique. Temperature, humidity, wind speed, wind direction, rainfall, incoming and reflected shortwave radiation and net radiation are measured above the canopy.

Soil heat fluxes are measured and soil moisture content is gathered using time domain reflectometry.

\subsection*{Purpose \& Scientific Goals}

\textbf{Purpose}

The primary purpose of the Howard Springs Flux Station is to understand the effects of fire on heat, moisture and carbon dioxide fluxes in Australia's tropical savannas. Other aims include:

\begin{itemize}
  \item to examine the water and carbon exchanges of tropical savannas
  \item understand the process of carbon cycling and storage in tropical savannas
  \item provide longer term measurements for future projects.
\end{itemize}

\subsection*{Site Photographs}

\begin{figure}[H]
\centering
  \includegraphics[width=0.8\textwidth,height=0.4\textheight,keepaspectratio]{/home/sanjays/et97_scratch2/oldscratch/Ozflux_data_full/ozflux_mirror/www.ozflux.org.au/monitoringsites/howardsprings/images/howard1.JPG}
  \caption*{Photo 1: A moderate intensity fire through our eddy covariance site}
\end{figure}

\begin{figure}[H]
\centering
  \includegraphics[width=0.8\textwidth,height=0.4\textheight,keepaspectratio]{/home/sanjays/et97_scratch2/oldscratch/Ozflux_data_full/ozflux_mirror/www.ozflux.org.au/monitoringsites/howardsprings/images/howard2.JPG}
  \caption*{Photo 2: Field Visit to the flux  tower, OzFlux Conference 2008}
\end{figure}

\begin{figure}[H]
\centering
  \includegraphics[width=0.8\textwidth,height=0.4\textheight,keepaspectratio]{/home/sanjays/et97_scratch2/oldscratch/Ozflux_data_full/ozflux_mirror/www.ozflux.org.au/monitoringsites/howardsprings/images/howard3.JPG}
  \caption*{Photo 3: Early wet season 2010}
\end{figure}

\begin{figure}[H]
\centering
  \includegraphics[width=0.8\textwidth,height=0.4\textheight,keepaspectratio]{/home/sanjays/et97_scratch2/oldscratch/Ozflux_data_full/ozflux_mirror/www.ozflux.org.au/monitoringsites/howardsprings/images/howard4.JPG}
  \caption*{Photo 4: Mid dry season 2010}
\end{figure}

\begin{figure}[H]
\centering
  \includegraphics[width=0.8\textwidth,height=0.4\textheight,keepaspectratio]{/home/sanjays/et97_scratch2/oldscratch/Ozflux_data_full/ozflux_mirror/www.ozflux.org.au/monitoringsites/howardsprings/images/howard5.JPG}
  \caption*{Photo 5: Solar power supply at  Howard Springs}
\end{figure}

\subsection*{Measurement Instrumentation}

\begin{longtable}{p{4.5cm}p{3.5cm}p{3.5cm}p{3.5cm}}
\toprule
\textbf{Instrument Type} & \textbf{Make} & \textbf{Model} & \textbf{Situation} \\
\midrule
\endhead
\bottomrule
\endfoot
Open path  CO2, H2O & LI-COR & LI-7500 & 23m \\
Sonic anemometer                 - wind velocities (u,v,w)                  - sonic   temperature & Campbell Scientific & CSAT-3 & 23m \\
Solar  Radiation (incoming) & Kipp  and Zonen & CM-7B & 23m \\
Solar  Radiation (reflected) & Kipp  and Zonen & CG-2 & 23m \\
Atmospheric  Pressure & Li-Cor & LI-7500 & 23m \\
Soil Heat Flux (2 replicates) & Campbell Scientific & HFT3 & Ground,  -15 cm \\
Soil moisture & Campbell Scientific & CS616 & 10-100cm \\
Temperature and Relative Humidity & Vaisala & HMP45A & 2m, 23m \\
Rain  Gauge & Campbell Scientific & TB3 & Ground \\
Data Loggers & Campbell Scientific & CR-3000 &  \\
Power Supply                 -                  12V DC EC flux station & Kyocera & Solar  panels (3) (BP) & Ground \\
\end{longtable}


\subsection*{Tower Global Metadata}

\begin{longtable}{p{4.5cm}p{10.5cm}}
\toprule
\textbf{Attribute} & \textbf{Value} \\
\midrule
\endhead
\bottomrule
\endfoot
\textbf{PI} & Lindsay Hutley and Jason Beringer \\
\textbf{Institution} & Charles Darwin University and University of Western Australia \\
\textbf{Latitude} & -12.4952 \\
\textbf{Longitude} & 131.15005 \\
\textbf{Elevation} & N/A \\
\textbf{Soil Type} & red kandasol \\
\textbf{Vegetation} & Woody savanna \\
\textbf{Temporal Coverage} & 2002-01-01 06:00:00 to 2025-07-08 10:30:00 \\
\textbf{Data Link} & \url{http://data.ozflux.org.au/} \\
\end{longtable}

\newpage

\section{Site Profile: Litchfield}

\subsection*{Physical Configuration \& Soil Sensor Depths}

\begin{figure}[H]
\centering
  \includegraphics[width=0.85\textwidth,height=0.45\textheight,keepaspectratio]{/home/sanjays/et97_scratch2/oldscratch/Ozflux_data_full/L6/soil_depth_plots/FluxTower_Diagrams/Litchfield/Litchfield_diagram.png}
  \caption*{Soil sensor depth profile and tower instrument heights.}
\end{figure}

\newpage

\subsection*{Site Description}

\textbf{Site Description}

The Litchfield flux station is located in the Litchfield National Park in the
          Northern Territory (GPS coordinates: -13.1790, 130.7945).

The station is a research and monitoring site representative of high rainfall, frequently burnt tropical savanna.
            The site is a 5 km x 5 km block of relatively uniform open-forest savanna inside the park about 80 km south of Darwin.

Measurements of carbon sequestration and remote sensing properties of vegetation through time are being achieved via
            a 40 m guyed tower equipped with instrumentation capable of directly measuring CO2, water use and surface
            energy properties (energy balance, reflectance).

\subsection*{Purpose \& Scientific Goals}

\textbf{Purpose}

The Litchfield flux station will provide nationally consistent observations of vegetation dynamics, faunal biodiversity, micrometeorology (climate,
            radiation, fluxes of carbon and water), hydrology and biogeochemistry to examine the impacts of fire regime, climate on carbon stocks
            and GHG emissions, and impacts on habitat quality via ongoing monitoring of vegetation structure and fauna. A wide range of ground based
            observations of vegetation structure and floristics is planned and all will link to remote sensing of fire and vegetation change over time.
            Co-location of Flux observations and remote sensing systems will be co-located to provide TERN with a ground-/air-/space based remote
            sensing observation stream. This will enable the development of tools describing fire occurrence, severity and associated greenhouse
            gas emissions, evapotranspiration and carbon sequestration

On-going fauna and flora monitoring will be undertaken at 5 sites within the SSS footprint that are component sites of the '3 Parks'
            data base, a long-term (15 years+) monitoring program across the 3 major national parks of the NT (Kakadu, Nitmiluk and Litchfield
            National Parks) that is examining the effects of fire frequency on vegetation structure and woody growth rates. Fire management,
            both protective and experimental burning regimes will be implemented across the SSS to assess impacts of fire on flora, fauna and
            biogeochemical cycles in savanna.

The tower will provide long-term measurements as part of the Ozflux network.

Key research questions

\begin{itemize}
  \item What are the impacts of prevailing fire regimes (primarily frequency, but also intensity, extent, heterogeneity) on vegetation
                structure and composition, habitat quality, fragmentation and vertebrate faunal biodiversity?
  \item How does vegetation structure, climate drivers and fire regime influence savanna carbon sequestration rate?
  \item How can fire management contribute to greenhouse gas abatement and carbon sequestration in savanna ecosystems?
  \item What are the impact of climate change on fire regimes and subsequent feedbacks to savanna carbon and water cycles?
\end{itemize}

\subsection*{Site Photographs}

\begin{figure}[H]
\centering
  \includegraphics[width=0.8\textwidth,height=0.4\textheight,keepaspectratio]{/home/sanjays/et97_scratch2/oldscratch/Ozflux_data_full/ozflux_mirror/www.ozflux.org.au/monitoringsites/litchfield/images/shed.jpg}
  \caption*{Photo 1: The  beginnings of the Litchfield flux station...}
\end{figure}

\begin{figure}[H]
\centering
  \includegraphics[width=0.8\textwidth,height=0.4\textheight,keepaspectratio]{/home/sanjays/et97_scratch2/oldscratch/Ozflux_data_full/ozflux_mirror/www.ozflux.org.au/monitoringsites/litchfield/images/crane.jpg}
  \caption*{Photo 2: The first part of the mast being erected}
\end{figure}

\begin{figure}[H]
\centering
  \includegraphics[width=0.8\textwidth,height=0.4\textheight,keepaspectratio]{/home/sanjays/et97_scratch2/oldscratch/Ozflux_data_full/ozflux_mirror/www.ozflux.org.au/monitoringsites/litchfield/images/baseFooting.jpg}
  \caption*{Photo 3: The footing of the tower base}
\end{figure}

\begin{figure}[H]
\centering
  \includegraphics[width=0.8\textwidth,height=0.4\textheight,keepaspectratio]{/home/sanjays/et97_scratch2/oldscratch/Ozflux_data_full/ozflux_mirror/www.ozflux.org.au/monitoringsites/litchfield/images/joiningTheTower.jpg}
  \caption*{Photo 4: The tower gets joined together}
\end{figure}

\begin{figure}[H]
\centering
  \includegraphics[width=0.8\textwidth,height=0.4\textheight,keepaspectratio]{/home/sanjays/et97_scratch2/oldscratch/Ozflux_data_full/ozflux_mirror/www.ozflux.org.au/monitoringsites/litchfield/images/onTheMast.jpg}
  \caption*{Photo 5: Climbing up}
\end{figure}

\begin{figure}[H]
\centering
  \includegraphics[width=0.8\textwidth,height=0.4\textheight,keepaspectratio]{/home/sanjays/et97_scratch2/oldscratch/Ozflux_data_full/ozflux_mirror/www.ozflux.org.au/monitoringsites/litchfield/images/lookingDown.jpg}
  \caption*{Photo 6: Looking down the mast}
\end{figure}

\begin{figure}[H]
\centering
  \includegraphics[width=0.8\textwidth,height=0.4\textheight,keepaspectratio]{/home/sanjays/et97_scratch2/oldscratch/Ozflux_data_full/ozflux_mirror/www.ozflux.org.au/monitoringsites/litchfield/images/groundAndCanopy.jpg}
  \caption*{Photo 7: Looking out on to the landcover from the mast}
\end{figure}

\begin{figure}[H]
\centering
  \includegraphics[width=0.8\textwidth,height=0.4\textheight,keepaspectratio]{/home/sanjays/et97_scratch2/oldscratch/Ozflux_data_full/ozflux_mirror/www.ozflux.org.au/monitoringsites/litchfield/images/canopyAndHorizon.jpg}
  \caption*{Photo 8: Smoke on the horizon}
\end{figure}

\begin{figure}[H]
\centering
  \includegraphics[width=0.8\textwidth,height=0.4\textheight,keepaspectratio]{/home/sanjays/et97_scratch2/oldscratch/Ozflux_data_full/ozflux_mirror/www.ozflux.org.au/monitoringsites/litchfield/images/groundAndCanopyDuringProtectiveBurn.jpg}
  \caption*{Photo 9: Looking out during a protective burn}
\end{figure}

\begin{figure}[H]
\centering
  \includegraphics[width=0.8\textwidth,height=0.4\textheight,keepaspectratio]{/home/sanjays/et97_scratch2/oldscratch/Ozflux_data_full/ozflux_mirror/www.ozflux.org.au/monitoringsites/litchfield/images/completedTower.jpg}
  \caption*{Photo 10: The completed tower}
\end{figure}

\subsection*{Measurement Instrumentation}

\begin{longtable}{p{4.5cm}p{3.5cm}p{3.5cm}p{3.5cm}}
\toprule
\textbf{Instrument Type} & \textbf{Make} & \textbf{Model} & \textbf{Location} \\
\midrule
\endhead
\bottomrule
\endfoot
Open path         CO2, H2O & Li-Cor & Li-7500A & 31m \\
Sonic anemometer & Campbell Scientific & CSAT-3 & 31m \\
Short and long wave radiation components, up and down & Kipp and Zonen & CNR4 & 31m \\
Net Radiometer & Kipp and Zonen & NRLite & 31m \\
Relative humidity and temperature & Vaisala & HMP155 & 33m             15m \\
Rainfall & Campbell Scientific & CS700L & 1m \\
Soil Heat Flux Plates & Hukseflux & HFP01 & 8cm (x2) \\
Soil temperature & Campbell Scientific & TCAV & 2cm       6cm \\
Soil moisture & Campbell Scientific & CS-616 & 5cm       20cm       50cm       100cm       140cm \\
Soil moisture and temperature & Campbell Scientific & CS-650 & 5cm       20cm       50cm       100cm       140cm \\
Phenocam & Wingscape & WCT-00122 & nadir, 30m       oblique (30$^\circ$), 30m (x2) \\
Web Camera & Vivotek & IP8362 & 35m \\
PAR Quantum sensor & LI-COR & LI-190 & facing up, 31m     facing down, 31m \\
4-channel light sensor & Skye Instruments & SKR 1850 & facing up, 31m     facing down, 31m \\
Total/direct/diffuse radiometer and sun hour sensor & Delta-T Devices & SPN1 & 31m \\
Acoustic sensor & Song Meter & SM2+ & 2m \\
\end{longtable}


\subsection*{Tower Global Metadata}

\begin{longtable}{p{4.5cm}p{10.5cm}}
\toprule
\textbf{Attribute} & \textbf{Value} \\
\midrule
\endhead
\bottomrule
\endfoot
\textbf{PI} & Lindsay Hutley and Jason Beringer \\
\textbf{Institution} & Charles Darwin University and University of Western Australia \\
\textbf{Latitude} & -13.17904178 \\
\textbf{Longitude} & 130.7945459 \\
\textbf{Elevation} & N/A \\
\textbf{Soil Type} & red kandasol \\
\textbf{Vegetation} & Woody savanna \\
\textbf{Temporal Coverage} & 2015-06-23 13:00:00 to 2025-07-08 10:00:00 \\
\textbf{Data Link} & \url{http://data.ozflux.org.au/} \\
\end{longtable}

\newpage

\section{Site Profile: Loxton}

\subsection*{Physical Configuration \& Soil Sensor Depths}

\begin{figure}[H]
\centering
  \includegraphics[width=0.85\textwidth,height=0.45\textheight,keepaspectratio]{/home/sanjays/et97_scratch2/oldscratch/Ozflux_data_full/L6/soil_depth_plots/FluxTower_Diagrams/Loxton/Loxton_diagram.png}
  \caption*{Soil sensor depth profile and tower instrument heights.}
\end{figure}

\newpage

\subsection*{Tower Global Metadata}

\begin{longtable}{p{4.5cm}p{10.5cm}}
\toprule
\textbf{Attribute} & \textbf{Value} \\
\midrule
\endhead
\bottomrule
\endfoot
\textbf{PI} & Robert Stevens and Cacilia Ewenz \\
\textbf{Institution} & South Australian Research and Development Institute/Flinders University/Airborne Research Australia \\
\textbf{Latitude} & -34.47035 \\
\textbf{Longitude} & 140.65512 \\
\textbf{Elevation} & N/A \\
\textbf{Soil Type} & loamy sand (0 - 50cm), sandy loam (50 - 85cm), sandy clay loam (85 - 170cm) \\
\textbf{Vegetation} & Almond orchard \\
\textbf{Temporal Coverage} & 2008-08-19 17:00:00 to 2009-06-09 10:30:00 \\
\textbf{Data Link} & \url{http://data.ozflux.org.au/} \\
\end{longtable}

\newpage

\section{Site Profile: MyallValeA}

\subsection*{Physical Configuration \& Soil Sensor Depths}

\begin{figure}[H]
\centering
  \includegraphics[width=0.85\textwidth,height=0.45\textheight,keepaspectratio]{/home/sanjays/et97_scratch2/oldscratch/Ozflux_data_full/L6/soil_depth_plots/FluxTower_Diagrams/MyallValeA/MyallValeA_diagram.png}
  \caption*{Soil sensor depth profile and tower instrument heights.}
\end{figure}

\newpage

\subsection*{Tower Global Metadata}

\begin{longtable}{p{4.5cm}p{10.5cm}}
\toprule
\textbf{Attribute} & \textbf{Value} \\
\midrule
\endhead
\bottomrule
\endfoot
\textbf{PI} & Ben Macdonald \\
\textbf{Institution} & TERN Landscapes/CSIRO \\
\textbf{Latitude} & -30.543 \\
\textbf{Longitude} & 150.02 \\
\textbf{Elevation} & N/A \\
\textbf{Soil Type} &  \\
\textbf{Vegetation} &  \\
\textbf{Temporal Coverage} & 2023-11-29 10:00:00 to 2025-01-17 00:00:00 \\
\textbf{Data Link} & \url{http://data.ozflux.org.au/} \\
\end{longtable}

\newpage

\section{Site Profile: MyallValeB}

\subsection*{Physical Configuration \& Soil Sensor Depths}

\begin{figure}[H]
\centering
  \includegraphics[width=0.85\textwidth,height=0.45\textheight,keepaspectratio]{/home/sanjays/et97_scratch2/oldscratch/Ozflux_data_full/L6/soil_depth_plots/FluxTower_Diagrams/MyallValeB/MyallValeB_diagram.png}
  \caption*{Soil sensor depth profile and tower instrument heights.}
\end{figure}

\newpage

\subsection*{Tower Global Metadata}

\begin{longtable}{p{4.5cm}p{10.5cm}}
\toprule
\textbf{Attribute} & \textbf{Value} \\
\midrule
\endhead
\bottomrule
\endfoot
\textbf{PI} & Ben Macdonald \\
\textbf{Institution} & TERN Landscapes/CSIRO \\
\textbf{Latitude} & -30.697 \\
\textbf{Longitude} & 150.09 \\
\textbf{Elevation} & N/A \\
\textbf{Soil Type} &  \\
\textbf{Vegetation} &  \\
\textbf{Temporal Coverage} & 2023-11-28 13:30:00 to 2025-01-17 00:00:00 \\
\textbf{Data Link} & \url{http://data.ozflux.org.au/} \\
\end{longtable}

\newpage

\section{Site Profile: Otway}

\subsection*{Physical Configuration \& Soil Sensor Depths}

\begin{figure}[H]
\centering
  \includegraphics[width=0.85\textwidth,height=0.45\textheight,keepaspectratio]{/home/sanjays/et97_scratch2/oldscratch/Ozflux_data_full/L6/soil_depth_plots/FluxTower_Diagrams/Otway/Otway_diagram.png}
  \caption*{Soil sensor depth profile and tower instrument heights.}
\end{figure}

\newpage

\subsection*{Site Description}

\textbf{Site Description}

The Otway flux station was located at Narrinda South in south west Victoria, Australia (GPS coordinates: -38.5323, 142.8168; elevation: 54m).

The pasture was grazed by dairy cattle with average grass height of 0.1m.

Annual average rainfall at the site was around 800mm and was only moderately seasonal. Mean daily temperature ranged from 25$^\circ$C in February to 12$^\circ$C in July.

The flux station was situated on a 10m tower. Fluxes of heat, water vapour and carbon dioxide were measured using the open-path eddy covariance technique. Supplementary measurements included temperature, humidity, rainfall, total solar, photosynthetically active radiation (PAR) and net radiation. Soil temperature and heat flux were also measured.

\subsection*{Purpose \& Scientific Goals}

\textbf{Purpose}

The purpose of the Otway flux station was:

\begin{itemize}
  \item to measure exchanges of carbon dioxide, water vapour and energy between the soil and the atmosphere using micrometeorological techniques.
  \item to participate in the atmospheric monitoring strategy, which was developed by CO2CRC researchers in CSIRO, and is based on the effects of hypothetical storage leaks (point and diffuse) simulated by atmospheric dispersion models, over several distance scales at the Otway site.
\end{itemize}

The Otway Project included injection and geological storage of up to 100,000 tonnes of carbon dioxide together with the most extensive monitoring and verification program yet undertaken at a geosequestration site.

\subsection*{Site Photographs}

\begin{figure}[H]
\centering
  \includegraphics[width=0.8\textwidth,height=0.4\textheight,keepaspectratio]{/home/sanjays/et97_scratch2/oldscratch/Ozflux_data_full/ozflux_mirror/www.ozflux.org.au/monitoringsites/otway/images/openpath.png}
  \caption*{Photo 1: Gill HS Sonic, open path IRGA and air intake for closed, heated open path   IRGA}
\end{figure}

\begin{figure}[H]
\centering
  \includegraphics[width=0.8\textwidth,height=0.4\textheight,keepaspectratio]{/home/sanjays/et97_scratch2/oldscratch/Ozflux_data_full/ozflux_mirror/www.ozflux.org.au/monitoringsites/otway/images/closed_openpath.png}
  \caption*{Photo 2: closed open path IRGA}
\end{figure}

\begin{figure}[H]
\centering
  \includegraphics[width=0.8\textwidth,height=0.4\textheight,keepaspectratio]{/home/sanjays/et97_scratch2/oldscratch/Ozflux_data_full/ozflux_mirror/www.ozflux.org.au/monitoringsites/otway/images/radThwind.png}
  \caption*{Photo 3: radiatian and temperature/humidity measurements as well as 2D   sonic}
\end{figure}

\subsection*{Measurement Instrumentation}

\begin{longtable}{p{4.5cm}p{3.5cm}p{3.5cm}p{3.5cm}}
\toprule
\textbf{Instrument Type} & \textbf{Make} & \textbf{Model} & \textbf{Situation} \\
\midrule
\endhead
\bottomrule
\endfoot
IRGA Open path        - density CO2, H2O        - atmospheric pressure & LI-COR & LI-7500 & 4.5m \\
3D sonic anemometer       - wind velocities (u,v,w)        - sonic   temperature & Gill & HS-3 & 4.5m \\
2D sonic anemometer x2 - wind velocities (u,v) & Gill & WindSonic & 0.3m located 1m from tower, 9.3m \\
Air temperature, Humidity x2       - T, Rh & Vaisala & HMP-45C & 0.2m located 1m from tower, 8.8m \\
Pyrgeometer Radiation        - Pyg & Kipp \& Zonen & CNR1 & 9m \\
Pyrgeometer Radiation       - Pyr & Kipp \& Zonen & CNR1 & 9m \\
Soil temperature x2        (averaging array) & Campbell & TCAV & -0.1m located 1m from tower \\
Soil water content x 2        (time domain reflectometer) & Campbell & CS616 & -0.1m located  1m from tower \\
Rainfall & Hydrological Services & TB3100 & 0.3m, located 1m  from base of tower \\
Data Loggers       - weather \& soil & Campbell & CR23X & 0m \\
PC  - EC fluxes & Advantech & ARK-8211 & 0m, located 250m from base of tower \\
Power Supply       - 24V DC EC flux station & Generic       - mains & Generic & 0m, located 250m from base of tower \\
\end{longtable}


\subsection*{Tower Global Metadata}

\begin{longtable}{p{4.5cm}p{10.5cm}}
\toprule
\textbf{Attribute} & \textbf{Value} \\
\midrule
\endhead
\bottomrule
\endfoot
\textbf{PI} & Eva van Gorsel \\
\textbf{Institution} & CSIRO Marine and Atmospheric Research \\
\textbf{Latitude} & -38.525 \\
\textbf{Longitude} & 142.81 \\
\textbf{Elevation} & N/A \\
\textbf{Soil Type} &  \\
\textbf{Vegetation} & Pasture, grazed \\
\textbf{Temporal Coverage} & 2007-08-11 00:00:00 to 2011-01-01 00:00:00 \\
\textbf{Data Link} & \url{http://data.ozflux.org.au/} \\
\end{longtable}

\newpage

\section{Site Profile: Red Dirt Melon Farm}

\subsection*{Physical Configuration \& Soil Sensor Depths}

\begin{figure}[H]
\centering
  \includegraphics[width=0.85\textwidth,height=0.45\textheight,keepaspectratio]{/home/sanjays/et97_scratch2/oldscratch/Ozflux_data_full/L6/soil_depth_plots/FluxTower_Diagrams/RedDirtMelonFarm/RedDirtMelonFarm_diagram.png}
  \caption*{Soil sensor depth profile and tower instrument heights.}
\end{figure}

\newpage

\subsection*{Tower Global Metadata}

\begin{longtable}{p{4.5cm}p{10.5cm}}
\toprule
\textbf{Attribute} & \textbf{Value} \\
\midrule
\endhead
\bottomrule
\endfoot
\textbf{PI} & Lindsay Hutley and Jason Beringer \\
\textbf{Institution} & Charles Darwin University and Monash University \\
\textbf{Latitude} & -14.563639 \\
\textbf{Longitude} & 132.477567 \\
\textbf{Elevation} & N/A \\
\textbf{Soil Type} & red kandasol \\
\textbf{Vegetation} & Woody savannah then cleared \\
\textbf{Temporal Coverage} & 2011-09-23 09:00:00 to 2013-07-21 23:30:00 \\
\textbf{Data Link} & \url{http://data.ozflux.org.au/} \\
\end{longtable}

\newpage

\section{Site Profile: Ridgefield}

\subsection*{Physical Configuration \& Soil Sensor Depths}

\begin{figure}[H]
\centering
  \includegraphics[width=0.85\textwidth,height=0.45\textheight,keepaspectratio]{/home/sanjays/et97_scratch2/oldscratch/Ozflux_data_full/L6/soil_depth_plots/FluxTower_Diagrams/Ridgefield/Ridgefield_diagram.png}
  \caption*{Soil sensor depth profile and tower instrument heights.}
\end{figure}

\newpage

\subsection*{Site Description}

\textbf{Site Description}

The Ridgefield flux station is located in at the UWA future farm in the wheatbelt of Western Australia about 200km southeast from Perth,
            Western Australia (GPS coordinates: -32.5061, 116.9668; elevation: 330m).  The Future Farm 2050 Project is based on UWA Farm Ridgefield. The Project's
            vision is to imagine the best-practice farm of 2050 and build and manage it now
            (http://www.ioa.uwa.edu.au/future-farm-2050).
            The site is established in a typical wheat field (a dominant managed land use type in Western Australia) and in an area that has been subject to extensive clearing.  We will contrast this with a native vegetated system and explore the past, present and future GHG budgets of these systems.

The flux tower site is located within an area of dryland agriculture. The surrounding area is dominated by broadacre farming practices. The vegetation cover is predominantly pasture. Elevation of the site is close to 330m and mean annual precipitation at a nearby Bureau of Meteorology site at Pingelly measures 446mm.

Maximum temperatures range from 15.3$^\circ$C (in July) to 31.9$^\circ$C (in Jan), while minimum temperatures range from 5.5$^\circ$C (in July) to 15.5$^\circ$C (in Jan).

The instrument mast is 10 meters tall. Heat, water vapour, carbon dioxide and methane measurements are taken using the open-path eddy flux technique. Temperature, humidity, wind speed, wind direction, rainfall and net radiation are measured. Soil heat fluxes are measured and soil moisture content are collected.

\subsection*{Purpose \& Scientific Goals}

\textbf{Purpose}

The purpose of the Ridgefield Flux Station is to:

\begin{itemize}
  \item Monitor and determine the balance of environmental demands for water yields, agricultural productivity, GHG budgets and biodiversity within a catchment landscape.
  \item Provide information to establish a modelling tool for GHG and water fluxes across various land use types, in order to benefit land management practices in the wheatbelt of Western Australia.
\end{itemize}

The tower is a crucial component of the UWA critical zone observatory which will be the site of a multidisciplinary approach to understanding the landscape dynamics.

\subsection*{Site Photographs}

\begin{figure}[H]
\centering
  \includegraphics[width=0.8\textwidth,height=0.4\textheight,keepaspectratio]{/home/sanjays/et97_scratch2/oldscratch/Ozflux_data_full/ozflux_mirror/www.ozflux.org.au/monitoringsites/ridgefield/images/rigdefield02.JPG}
  \caption*{Photo 1: International students at theInaugural World University Network Critical Zone Summer School}
\end{figure}

\begin{figure}[H]
\centering
  \includegraphics[width=0.8\textwidth,height=0.4\textheight,keepaspectratio]{/home/sanjays/et97_scratch2/oldscratch/Ozflux_data_full/ozflux_mirror/www.ozflux.org.au/monitoringsites/ridgefield/images/rigdefield03.JPG}
  \caption*{Photo 2: Completed tower}
\end{figure}

\begin{figure}[H]
\centering
  \includegraphics[width=0.8\textwidth,height=0.4\textheight,keepaspectratio]{/home/sanjays/et97_scratch2/oldscratch/Ozflux_data_full/ozflux_mirror/www.ozflux.org.au/monitoringsites/ridgefield/images/rigdefield01.JPG}
  \caption*{Photo 3: Flux instrumentation}
\end{figure}

\begin{figure}[H]
\centering
  \includegraphics[width=0.8\textwidth,height=0.4\textheight,keepaspectratio]{/home/sanjays/et97_scratch2/oldscratch/Ozflux_data_full/ozflux_mirror/www.ozflux.org.au/monitoringsites/ridgefield/images/rigdefield04.JPG}
  \caption*{Photo 4: Panorama of the site}
\end{figure}

\subsection*{Measurement Instrumentation}

\begin{longtable}{p{4.5cm}p{3.5cm}p{3.5cm}p{3.5cm}}
\toprule
\textbf{Instrument Type} & \textbf{Make} & \textbf{Model} & \textbf{Situation} \\
\midrule
\endhead
\bottomrule
\endfoot
Atmospheric Pressure & LI-COR & LI-7550A & 1.0m \\
Net Radiation & Kipp and Zonen & CNR4 & 8m \\
Data logger & Sutron & Xlite & Ground \\
Open Path CO2 \& H2O & LI-COR & LI-7550A \& LI-7500 & 2.5m \\
Open Path CH4 & LI-COR & LI-7700 & 2.5m \\
Rain gauge & Texas Electronics & CS702 & Ground \\
Soil heat flux (three replicates) & Hukseflux self calibrating probes & HFP01 & 0.1m \\
Soil temperature & Stevens & Hydra Probe II & 0.05, 0.10, 0.20m \\
Soil moisture & Stevens & Hydra Probe II & 0.05, 0.10, 0.20m \\
PAR & LI-COR & LI-190 &  \\
Sonic anemometer                 - wind velocities (u,v,w)                 - sonic temperature & Campbell Scientific & CSAT3 & 2.5m \\
Temperature and Relative Humidity & Vaisala & HMP155 & 2.5m \\
\end{longtable}


\subsection*{Tower Global Metadata}

\begin{longtable}{p{4.5cm}p{10.5cm}}
\toprule
\textbf{Attribute} & \textbf{Value} \\
\midrule
\endhead
\bottomrule
\endfoot
\textbf{PI} & Jason Beringer \\
\textbf{Institution} & University of Western Australia \\
\textbf{Latitude} & -32.506102 \\
\textbf{Longitude} & 116.966827 \\
\textbf{Elevation} & N/A \\
\textbf{Soil Type} &  \\
\textbf{Vegetation} & Wheat \\
\textbf{Temporal Coverage} & 2015-03-29 05:00:00 to 2025-07-08 08:30:00 \\
\textbf{Data Link} & \url{http://data.ozflux.org.au/} \\
\end{longtable}

\newpage

\section{Site Profile: Riggs Creek}

\subsection*{Physical Configuration \& Soil Sensor Depths}

\begin{figure}[H]
\centering
  \includegraphics[width=0.85\textwidth,height=0.45\textheight,keepaspectratio]{/home/sanjays/et97_scratch2/oldscratch/Ozflux_data_full/L6/soil_depth_plots/FluxTower_Diagrams/RiggsCreek/RiggsCreek_diagram.png}
  \caption*{Soil sensor depth profile and tower instrument heights.}
\end{figure}

\newpage

\subsection*{Site Description}

\textbf{Site Description}

The Riggs  Creek flux station was located in Goulburn-Broken catchment in North-Eastern  Victoria (GPS coordinates: -36.6499, 145.5760).

The flux tower site was located within an area of dryland agriculture. The surrounding area is dominated  by broadacre farming practices. The vegetation cover is predominantly pasture. Elevation of the site is close to 152m and mean annual  precipitation at a nearby Bereau of Meteorology site measures 650mm.

Maximum  temperatures range from 12.3$^\circ$C (in July) to 29.7$^\circ$C (in  February), while minimum temperatures range from 10.4$^\circ$C (in July) to  26.8$^\circ$C (in February).

The instrument mast was 4 meters  tall. Heat, water vapour and carbon dioxide measurements were taken using the  open-path eddy flux technique. Temperature, humidity, wind speed, wind  direction, rainfall and net radiation were measured. Soil heat fluxes were  measured and soil moisture content was gathered using time domain  reflectometry.

\subsection*{Purpose \& Scientific Goals}

\textbf{Purpose}

The purpose of the Riggs  Creek Flux Station was to:

\begin{itemize}
  \item Monitor  and determine the balance of environmental demands for water yields,  agricultural productivity, carbon storage and biodiversity within a catchment  landscape.
  \item Provide  information for other projects (i.e. The Carbon Project ) which seeks to establish a  modelling tool for carbon and water fluxes across various land use types, in  order to benefit land management practices in the Goulburn-Broken catchment  region.
\end{itemize}

\subsection*{Site Photographs}

\begin{figure}[H]
\centering
  \includegraphics[width=0.8\textwidth,height=0.4\textheight,keepaspectratio]{/home/sanjays/et97_scratch2/oldscratch/Ozflux_data_full/ozflux_mirror/www.ozflux.org.au/monitoringsites/riggscreek/images/Riggs1.jpg}
  \caption*{Photo 1: Wiring  the tower, November 2010}
\end{figure}

\begin{figure}[H]
\centering
  \includegraphics[width=0.8\textwidth,height=0.4\textheight,keepaspectratio]{/home/sanjays/et97_scratch2/oldscratch/Ozflux_data_full/ozflux_mirror/www.ozflux.org.au/monitoringsites/riggscreek/images/Riggs4.jpg}
  \caption*{Photo 2: final touches to the tower, November 2010}
\end{figure}

\begin{figure}[H]
\centering
  \includegraphics[width=0.8\textwidth,height=0.4\textheight,keepaspectratio]{/home/sanjays/et97_scratch2/oldscratch/Ozflux_data_full/ozflux_mirror/www.ozflux.org.au/monitoringsites/riggscreek/images/Riggs2.JPG}
  \caption*{Photo 3: The  completed tower}
\end{figure}

\begin{figure}[H]
\centering
  \includegraphics[width=0.8\textwidth,height=0.4\textheight,keepaspectratio]{/home/sanjays/et97_scratch2/oldscratch/Ozflux_data_full/ozflux_mirror/www.ozflux.org.au/monitoringsites/riggscreek/images/Riggs3.jpg}
  \caption*{Photo 4: Team  photograph after the successful construction of the tower, November 2010}
\end{figure}

\subsection*{Measurement Instrumentation}

\begin{longtable}{p{4.5cm}p{3.5cm}p{3.5cm}p{3.5cm}}
\toprule
\textbf{Instrument Type} & \textbf{Make} & \textbf{Model} & \textbf{Location} \\
\midrule
\endhead
\bottomrule
\endfoot
Atmospheric  Pressure & Vaisala & CS106 & 2.5m \\
Net  Radiation & Campbell Scientific & CNR4 & 4m \\
Data  logger & Campbell Scientific & CR3000 & Ground \\
Open  Path CO2 H2O & Li-Cor & Li-7500 & 2.5m \\
Rain gauge & Campbell Scientific & CS702 & Ground \\
Soil  heat flux (two replicates) & Hukseflux & HFP01 & 0.1m \\
Soil  temperature & Campbell Scientific & TCAV-L & 0.08m \\
Sonic anemometer                 - wind velocities (u,v,w)                  - sonic temperature & Campbell Scientific & CSAT3 & 2.5m \\
Temperature  and Relative Humidity  (two replicates) & Vaisala & HMP45AC & 2.5m \\
Soil  Moisture (6 replicates) & Campbell Scientific & CS616 & 0.1,  0.2, 0.4, 0.8, 1.2, 2.0m \\
\end{longtable}


\subsection*{Tower Global Metadata}

\begin{longtable}{p{4.5cm}p{10.5cm}}
\toprule
\textbf{Attribute} & \textbf{Value} \\
\midrule
\endhead
\bottomrule
\endfoot
\textbf{PI} & Jason Beringer \\
\textbf{Institution} & University of Western Australia \\
\textbf{Latitude} & -36.656 \\
\textbf{Longitude} & 145.576 \\
\textbf{Elevation} & N/A \\
\textbf{Soil Type} & Unknown \\
\textbf{Vegetation} & Pasture \\
\textbf{Temporal Coverage} & 2011-01-01 00:30:00 to 2017-07-12 17:30:00 \\
\textbf{Data Link} & \url{http://data.ozflux.org.au/} \\
\end{longtable}

\newpage

\section{Site Profile: Robson Creek}

\subsection*{Physical Configuration \& Soil Sensor Depths}

\begin{figure}[H]
\centering
  \includegraphics[width=0.85\textwidth,height=0.45\textheight,keepaspectratio]{/home/sanjays/et97_scratch2/oldscratch/Ozflux_data_full/L6/soil_depth_plots/FluxTower_Diagrams/RobsonCreek/RobsonCreek_diagram.png}
  \caption*{Soil sensor depth profile and tower instrument heights.}
\end{figure}

\newpage

\subsection*{Site Description}

\textbf{Site Description}

The Robson Creek tower is located in tropical rainforest, approximately 30 km northwest of Atherton in Far North Queensland (GPS coordinates: -17.1175, 145.6301; elevation: 710 m).

It lies on the western slopes of the Lamb Range in Danbulla National Park, within the Wet Tropics World Heritage Area (WTWHA).

The tower is located to the NW of a 25 Ha census plot established by CSIRO in 2012.

Vegetation type: the forest is mapped as Regional Ecosystem (RE) 7.3.36a, complex mesophyll vine forest (Queensland Government 2006). The forest type changes to RE 7.12.16a, simple to complex notophyll vine forest, with increasing altitude to the north of the 25 Ha plot.

In structure the forest is very tall to extremely tall closed forest with canopy heights ranging from 23 to 44 m.

The tower is located at the southern base of the Lamb Range which rises to 1276 m ASL.

The landform of the 25Ha plot which is in the dominant wind direction from the tower is moderately inclined with a low relief, a 30 m high ridge running north/south through the middle of the plot and a 40 m high ridge running north/south on the eastern edge of the plot.

Three permanent creeks flow through the 25 Ha plot, joining with Robson Creek which in turn meets the Tinaroo Dam approximately one kilometre south of the plot.

The climate of the area is considered seasonal with 61\% of the annual rainfall occurring in the months of January to March.

Based on the first 10 years of site rainfall data the MAP 2292 mm indicates the OzFlux site is slightly wetter than the nearby locations.

\subsection*{Purpose \& Scientific Goals}

\textbf{Purpose}

The purpose of the Robson Creek flux station is:

\begin{itemize}
  \item to measure exchanges of carbon dioxide, water vapour and energy between the tropical upland rainforest in Far North Queensland and the atmosphere using micrometeorological techniques
  \item to quantify the changes in carbon and water balances of an Australian tropical rainforest on a long term basis in the face of climate change
  \item to present the results from the study in real time to the public and inform the public on what these results mean.
\end{itemize}

This flux station provides primary site meteorological and radiometric data for Robson Creek, a reference calibration site (GEOTREES) with a 25 Ha Census vegetation plot.

\subsection*{Site Photographs}

\begin{figure}[H]
\centering
  \includegraphics[width=0.8\textwidth,height=0.4\textheight,keepaspectratio]{/home/sanjays/et97_scratch2/oldscratch/Ozflux_data_full/ozflux_mirror/www.ozflux.org.au/monitoringsites/robsoncreek/images/RobsonCreek_Tower1.jpg}
  \caption*{Photo 1: Robson Creek ‘small world’ view of the  instrument array from the tower top}
\end{figure}

\begin{figure}[H]
\centering
  \includegraphics[width=0.8\textwidth,height=0.4\textheight,keepaspectratio]{/home/sanjays/et97_scratch2/oldscratch/Ozflux_data_full/ozflux_mirror/www.ozflux.org.au/monitoringsites/robsoncreek/images/RobsonCreek_Tower2.jpg}
  \caption*{Photo 2: Robson Creek 40m tower view of the solar array from the tower base}
\end{figure}

\begin{figure}[H]
\centering
  \includegraphics[width=0.8\textwidth,height=0.4\textheight,keepaspectratio]{/home/sanjays/et97_scratch2/oldscratch/Ozflux_data_full/ozflux_mirror/www.ozflux.org.au/monitoringsites/robsoncreek/images/RobsonCreek_Tower3.jpg}
  \caption*{Photo 3: Robson Creek 40m tower view of the instrument array from the north}
\end{figure}

\subsection*{Measurement Instrumentation}

\begin{longtable}{p{4.5cm}p{3.5cm}p{3.5cm}p{3.5cm}}
\toprule
\textbf{Instrument Type} & \textbf{Make} & \textbf{Model} & \textbf{Situation} \\
\midrule
\endhead
\bottomrule
\endfoot
IRGA Open path                  - density CO2, H2O  LI-COR LI-7500RS 40m & LI-COR & LI-7500RS & 40m \\
IRGA Interface Unit                 - atmospheric pressure & LI-COR & LI-7550 & 40m \\
3D sonic anemometer                 - wind velocities (u,v,w)                  - sonic   temperature & Campbell Scientific & CSAT3B & 40m \\
Air temperature, Humidity                  - T, Rh & Vaisala & HMP155 & 40m \\
PAR Radiation                  - PPFD & LI-COR & LI-190SB & 40m \\
4 Component Radiation                - SW-d, LW-d, LW-u, SW-u & Hukseflux & NR01 & 40m \\
Net Radiation               - Rn & Kipp \& Zonen & NR Lite & 40m \\
Rainfall                 - Rain & Campbell Scientific & RIM8000 & 40m \\
Pyranometer & Apogee & SP-110 & 40m \\
Soil temperature                  (averaging array) & Campbell Scientific & TCAV & -7cm                 10m S of tower \\
Soil heat flux plates (2) & Hukseflux & HFP01 & -10cm                  10m S of tower \\
Soil water content                (time  domain reflectometer) & Campbell Scientific & CS616 & -6cm, -75cm,                 -150cm, -200cm                                  10m S of tower \\
Data Logger                 - EC fluxes & Campbell Scientific & CR1000X & 40m \\
CO2 sensor & Campbell Scientific & GMP343 & 40m \\
Phenocam - Field Camera 5MP                 - 30 minute images & Campbell Scientific & CCFC-1879 & 40m \\
Automatic weatherstation                 (T, Rh, Rain) & Vaisala & WXT530 & 40m \\
\end{longtable}


\subsection*{Tower Global Metadata}

\begin{longtable}{p{4.5cm}p{10.5cm}}
\toprule
\textbf{Attribute} & \textbf{Value} \\
\midrule
\endhead
\bottomrule
\endfoot
\textbf{PI} & Michael Liddell \\
\textbf{Institution} & James Cook University, Cairns \\
\textbf{Latitude} & -17.11746943 \\
\textbf{Longitude} & 145.6301375 \\
\textbf{Elevation} & N/A \\
\textbf{Soil Type} & Acidic, Dystrophic, Brown Dermosol \\
\textbf{Vegetation} & Regional Ecosystem (RE) 7.3.36a, complex mesophyll vine forest \\
\textbf{Temporal Coverage} & 2013-08-01 17:30:00 to 2025-07-23 08:30:00 \\
\textbf{Data Link} & \url{http://data.ozflux.org.au/} \\
\end{longtable}

\newpage

\section{Site Profile: SilverPlains}

\subsection*{Physical Configuration \& Soil Sensor Depths}

\begin{figure}[H]
\centering
  \includegraphics[width=0.85\textwidth,height=0.45\textheight,keepaspectratio]{/home/sanjays/et97_scratch2/oldscratch/Ozflux_data_full/L6/soil_depth_plots/FluxTower_Diagrams/SilverPlains/SilverPlains_diagram.png}
  \caption*{Soil sensor depth profile and tower instrument heights.}
\end{figure}

\newpage

\subsection*{Site Description}

\textbf{Site Description}

The Silver Plains flux tower is located at Interlaken, on the Tasmanian central plateau, Australia (GPS coordinates -42.15, 147.13, elevation 880m), on land owned and managed by the Tasmanian Land Conservancy.

The site is a natural grassy sedgeland with an average summer temperature of 16$^\circ$C, average winter temperature of 6$^\circ$C, and average annual rainfall of approximately 900mm. Soil at the site is peaty, being an organosol containing on average 8 kgCm-2 in the top 10cm, overlying black vertosol on Jurassic dolerite. The vegetation at the site is heavily grazed year round by a range of native vertebrate herbivores, including wallabies, pademelons and wombats, as well as by feral fallow deer, resulting in an extremely low vegetation stature of a few centimetres, with the exception of inflorescences which can extend up to 30cm above the ground.

The instrument tower is 3m tall, with carbon dioxide, water and heat fluxes being measured using an open-path eddy flux technique at 2.5m. Additional measurements include wind speed and direction, air temperature, humidity, incoming and outgoing short and longwave radiation, rainfall, soil temperature and soil water content.

\subsection*{Purpose \& Scientific Goals}

\textbf{Purpose}

The purpose of the Silver Plains flux tower is to:

\begin{itemize}
  \item Monitor exchanges of carbon dioxide, water vapour and energy in a high-altitude grassy peatland ecosystem on the Tasmanian Central Plateau.
  \item Quantify the carbon balance of the ecosystem, along with the key components of net ecosystem exchange, gross primary productivity and ecosystem respiration.
  \item Identify key environmental and climatic drivers of carbon, water and energy fluxes.
  \item Utilise measurements alongside manual chamber measurements and ancillary environmental variables to model ecosystem carbon dynamics under a future climate.
  \item to compare and contrast ecosystem function between unmanaged and irrigated pastures and native woodland vegetation.
\end{itemize}

\subsection*{Site Photographs}

\begin{figure}[H]
\centering
  \includegraphics[width=0.8\textwidth,height=0.4\textheight,keepaspectratio]{/home/sanjays/et97_scratch2/oldscratch/Ozflux_data_full/ozflux_mirror/www.ozflux.org.au/monitoringsites/silverplains/images/silverplainsECSystem.jpg}
  \caption*{Photo 1: The Silver Plains EC tower.}
\end{figure}

\begin{figure}[H]
\centering
  \includegraphics[width=0.8\textwidth,height=0.4\textheight,keepaspectratio]{/home/sanjays/et97_scratch2/oldscratch/Ozflux_data_full/ozflux_mirror/www.ozflux.org.au/monitoringsites/silverplains/images/silverplains.jpg}
  \caption*{Photo 2: Silver Plains}
\end{figure}

\subsection*{Measurement Instrumentation}

\begin{longtable}{p{4.5cm}p{3.5cm}p{3.5cm}p{3.5cm}}
\toprule
\textbf{Instrument Type} & \textbf{Make} & \textbf{Model} & \textbf{Situation} \\
\midrule
\endhead
\bottomrule
\endfoot
Open path  CO2, H2O & Campbell Scientific & EC150 & 2.5m \\
sonic anemometer & Campbell Scientific & CSAT3A & 2.5m \\
Sonic anemometer & Gill & Windsonic1 2D & 3m \\
Net Radiometer & Kipp and Zonen & CNR1 & 1.8m \\
Net Radiometer & Kipp and Zonen & NR-LITE2 & 1.8m \\
Quantum sensor & Apogee & SQ-110 PAR & facing up, 2.5mfacing down, 2.5m \\
RH+T & Campbell Scientific & EE181 & 2.5m \\
Thermistor & Betatherm & 100K6A1IA & 2m \\
Rainfall & HyQuest & TB4 & 1m \\
Soil Heat Flux & HukseFlux & HFP01 & 8cm (x3) \\
Soil temperature & Campbell Scientific & 107 & 30cm \\
Soil temperature & Campbell Scientific & TCAV & 1cm2cm4cm8cm \\
Soil moisture and temperature & Campbell Scientific & CS650 & 10cm (x4) 30cm (x2) \\
\end{longtable}


\subsection*{Tower Global Metadata}

\begin{longtable}{p{4.5cm}p{10.5cm}}
\toprule
\textbf{Attribute} & \textbf{Value} \\
\midrule
\endhead
\bottomrule
\endfoot
\textbf{PI} & Mark Hovenden \\
\textbf{Institution} & University of Tasmania \\
\textbf{Latitude} & -42.090556 \\
\textbf{Longitude} & 147.0875 \\
\textbf{Elevation} & N/A \\
\textbf{Soil Type} & Organosol on underlying dolerite geology \\
\textbf{Vegetation} & Species rich highland grassland/sedgeland \\
\textbf{Temporal Coverage} & 2019-12-05 10:00:00 to 2024-12-31 23:30:00 \\
\textbf{Data Link} & \url{http://data.ozflux.org.au/} \\
\end{longtable}

\newpage

\section{Site Profile: Snow Gum}

\subsection*{Physical Configuration \& Soil Sensor Depths}

\begin{figure}[H]
\centering
  \includegraphics[width=0.85\textwidth,height=0.45\textheight,keepaspectratio]{/home/sanjays/et97_scratch2/oldscratch/Ozflux_data_full/L6/soil_depth_plots/FluxTower_Diagrams/SnowGum/SnowGum_diagram.png}
  \caption*{Soil sensor depth profile and tower instrument heights.}
\end{figure}

\newpage

\subsection*{Tower Global Metadata}

\begin{longtable}{p{4.5cm}p{10.5cm}}
\toprule
\textbf{Attribute} & \textbf{Value} \\
\midrule
\endhead
\bottomrule
\endfoot
\textbf{PI} & Mark Hovenden \\
\textbf{Institution} & AMRF; ANU; UTAS \\
\textbf{Latitude} & -36.375649 \\
\textbf{Longitude} & 148.43585 \\
\textbf{Elevation} & N/A \\
\textbf{Soil Type} &  \\
\textbf{Vegetation} & Sub-alpine woodland \\
\textbf{Temporal Coverage} & 2023-12-27 15:30:00 to 2025-01-17 00:00:00 \\
\textbf{Data Link} & \url{http://data.ozflux.org.au/} \\
\end{longtable}

\newpage

\section{Site Profile: Sturt Plains}

\subsection*{Physical Configuration \& Soil Sensor Depths}

\begin{figure}[H]
\centering
  \includegraphics[width=0.85\textwidth,height=0.45\textheight,keepaspectratio]{/home/sanjays/et97_scratch2/oldscratch/Ozflux_data_full/L6/soil_depth_plots/FluxTower_Diagrams/SturtPlains/SturtPlains_diagram.png}
  \caption*{Soil sensor depth profile and tower instrument heights.}
\end{figure}

\newpage

\subsection*{Site Description}

\textbf{Site Description}

The Sturt  Plains flux station is located approximately 280km north of Tennant Creek,  Northern Territory (GPS coordinates: -17.1507, 133.3502).

The flux tower site is 280km north of Tennant Creek, amongst a low lying  plain dominated by \textit{Mitchell Grass (gen.}\textit{Astrebla)}\textit{.} Elevation of the site is close to 250m and mean annual precipitation at a nearby  Bureau of Meteorology site is 640mm.

Maximum temperatures range from 28.4$^\circ$C  (in June/ July) to 39.1$^\circ$C (in December), while minimum temperatures  range from 11.2$^\circ$C (in July) to 24.4$^\circ$C (in December).

The instrument mast is 5 meters  tall. Heat, water vapour and carbon dioxide measurements are taken using the  open-path eddy flux technique. Temperature, humidity, wind speed, wind  direction, rainfall and net radiation are measured. Soil heat fluxes are measured  and soil moisture content is gathered using time domain reflectometry.

Ancillary measurements taken at  the site include LAI, leaf-scale physiological properties (gas exchange, leaf  isotope ratios, N and chlorophyll concentrations), vegetation optical  properties and soil physical properties. Airborne based remote sensing (Lidar  and hyperspectral measurements) was carried out across the transect in  September 2008.

\subsection*{Purpose \& Scientific Goals}

\textbf{Purpose}

The  purpose of the Sturt Plains Flux Station is to:\textbf{}

\begin{itemize}
  \item Provide  information as part of a larger network of flux towers established along the  North Australian Tropical Transect (NATT) gradient, which extends \textasciitilde{}1000km  south from Darwin 12.5$^\circ$S.  The towers provide flux data across the savanna’s heterogeneous ‘top end’  ecosystems including open forest savanna, open savanna woodland and seasonally  inundated floodplains.
  \item Examine  spatial patterns and processes of land-surface-atmosphere exchanges (radiation,  heat, moisture, CO2 and other trace gasses) across scales from leaf to  landscape scales within Australian savannas.
  \item Determine  the climate and ecosystem characteristics (physical structure, species  composition, physiological function) that drive spatial and temporal variations  of carbon, water and energy fluxes from north Australian savanna.
  \item Determine  if fluxes of carbon, water vapour and heat over the various ecosystems as  derived from the various measurement techniques can be combined to form a  comprehensive and consistent estimate of the regional fluxes and budgets across  the landscape.
  \item Provide  long-term consistent observations of energy, carbon and water exchange between  the atmosphere and the ‘top end’ grassland.
\end{itemize}

\subsection*{Site Photographs}

\begin{figure}[H]
\centering
  \includegraphics[width=0.8\textwidth,height=0.4\textheight,keepaspectratio]{/home/sanjays/et97_scratch2/oldscratch/Ozflux_data_full/ozflux_mirror/www.ozflux.org.au/monitoringsites/sturtplains/images/sturt1.JPG}
  \caption*{Photo 1: The tower site shortly after it was assembled in August 2008}
\end{figure}

\begin{figure}[H]
\centering
  \includegraphics[width=0.8\textwidth,height=0.4\textheight,keepaspectratio]{/home/sanjays/et97_scratch2/oldscratch/Ozflux_data_full/ozflux_mirror/www.ozflux.org.au/monitoringsites/sturtplains/images/sturt2.JPG}
  \caption*{Photo 2: The tower site in April 2010}
\end{figure}

\begin{figure}[H]
\centering
  \includegraphics[width=0.8\textwidth,height=0.4\textheight,keepaspectratio]{/home/sanjays/et97_scratch2/oldscratch/Ozflux_data_full/ozflux_mirror/www.ozflux.org.au/monitoringsites/sturtplains/images/sturt3.JPG}
  \caption*{Photo 3: Variables beyond our control}
\end{figure}

\subsection*{Measurement Instrumentation}

\begin{longtable}{p{4.5cm}p{3.5cm}p{3.5cm}p{3.5cm}}
\toprule
\textbf{Instrument Type} & \textbf{Make} & \textbf{Model} & \textbf{Location} \\
\midrule
\endhead
\bottomrule
\endfoot
Albedometer & Kipp  \& Zonen & CM7B & 4.8m \\
Data  logger & Campbell Scientific & CR3000 & Ground \\
Open  Path CO2 and H2O & LiCor & Li-7500 & 4.8m \\
Net  Radiometer & Kipp  \& Zonen & N-R  Lite & 4.8m \\
Rain gauge & Campbell Scientific & CS702 & 4.8m \\
Soil  heat flux & Campbell Scientific & HFT3 & -0.15m \\
Soil  moisture (two replicates) & Campbell Scientific & CS616 & -0.05,  -0.50m \\
Sonic anemometer                 - wind velocities (u,v,w)                  - sonic temperature & Campbell Scientific & CSAT3 & 4.8m \\
Temperature  and Relative Humidity & Vaisala & HMP45AC & 4.8m \\
Net  Longwave Radiation & Kipp  \& Zonen & CG-2 & 4.8m \\
\end{longtable}


\subsection*{Tower Global Metadata}

\begin{longtable}{p{4.5cm}p{10.5cm}}
\toprule
\textbf{Attribute} & \textbf{Value} \\
\midrule
\endhead
\bottomrule
\endfoot
\textbf{PI} & Lindsay Hutley and Jason Beringer \\
\textbf{Institution} & Charles Darwin University and University of Western Australia \\
\textbf{Latitude} & -17.1507 \\
\textbf{Longitude} & 133.3502 \\
\textbf{Elevation} & N/A \\
\textbf{Soil Type} & red kandasol \\
\textbf{Vegetation} & Mitchell Grasslands \\
\textbf{Temporal Coverage} & 2008-08-28 11:00:00 to 2025-03-11 15:00:00 \\
\textbf{Data Link} & \url{http://data.ozflux.org.au/} \\
\end{longtable}

\newpage

\section{Site Profile: TiTreeEast}

\subsection*{Physical Configuration \& Soil Sensor Depths}

\begin{figure}[H]
\centering
  \includegraphics[width=0.85\textwidth,height=0.45\textheight,keepaspectratio]{/home/sanjays/et97_scratch2/oldscratch/Ozflux_data_full/L6/soil_depth_plots/FluxTower_Diagrams/TiTreeEast/TiTreeEast_diagram.png}
  \caption*{Soil sensor depth profile and tower instrument heights.}
\end{figure}

\newpage

\subsection*{Site Description}

\textbf{Site Description}

The Ti Tree East flux station is  located on Pine Hill cattle station in the Northern Territory (GPS coordinates: -22.2870, 133.6400).

The site is a mosaic of the primary  semi-arid biomes of central Australia: grassy mulga woodland and \textit{Corymbia}/\textit{Triodia} savanna.

The woodland is characterised by a mulga (\textit{Acacia aneura}) canopy, which is 4.85 m  tall on average.

Elevation of the site  is 553 m above sea level, and the terrain is flat.

Mean annual precipitation at the nearby (30 km  to the south) Bureau of Meteorology station is 305.9 mm but ranges between 100 mm  in 2009 to 750 mm in 2010.

Predominant  wind directions are from the southeast and east.

The soil is red sand overlying an 8 m deep  water table.

Pine Hill Station is a  functioning cattle station that has been in operation for longer than 50 years. However, the east side has not been stocked  in over three years.

The instrument mast is 10 m  tall. Fluxes of heat, water vapour and  carbon are measured using the open-path eddy covariance technique at 9.81 m.

Supplementary measurements above the canopy  include temperature and humidity (9.81 m), windspeed and wind direction (8.28 m),  downwelling and upwelling shortwave and longwave radiation (9.9 m).

Precipitation is monitored in the savanna  (2.5m). Supplementary measurements  within and below the canopy include barometric pressure (2 m).

Belowground soil measurements are made beneath \textit{Triodia}, mulga and grassy understorey  and include ground heat flux (0.08 m), soil temperature (0.02 m – 0.06 m) and  soil moisture (0 – 0.1 m, 0.1 – 0.3 m, 0.6 – 0.8 m and 1.0 – 1.2 m).

\subsection*{Purpose \& Scientific Goals}

\textbf{Purpose}

The  purpose of the Ti Tree  flux station is:

\begin{itemize}
  \item to measure the exchanges of carbon dioxide,  water vapour and energy in a semi-arid ecosystem with potential access to  groundwater
  \item to identify flux footprints associated with  contributions by mulga versus Corymbia savannas
  \item to compare water use efficiency, GPP and  ecosystem respiration between adjacent semi-arid ecosystems (Alice Springs  mulga)
  \item to identify relationships between groundwater,  soil moisture, rainfall and evapotranspiration
  \item to identify phenological trends and to relate  phenology to flux footprints and remote sensing of water and carbon fluxes.
\end{itemize}

\subsection*{Site Photographs}

\begin{figure}[H]
\centering
  \includegraphics[width=0.8\textwidth,height=0.4\textheight,keepaspectratio]{/home/sanjays/et97_scratch2/oldscratch/Ozflux_data_full/ozflux_mirror/www.ozflux.org.au/monitoringsites/titreeeast/images/metaTTE-photo1.jpg}
  \caption*{Photo 1: View from the top of the  tower (10m)}
\end{figure}

\begin{figure}[H]
\centering
  \includegraphics[width=0.8\textwidth,height=0.4\textheight,keepaspectratio]{/home/sanjays/et97_scratch2/oldscratch/Ozflux_data_full/ozflux_mirror/www.ozflux.org.au/monitoringsites/titreeeast/images/metaTTE-photo2.jpg}
  \caption*{Photo 2: View of Triodia/Corymbia savanna}
\end{figure}

\begin{figure}[H]
\centering
  \includegraphics[width=0.8\textwidth,height=0.4\textheight,keepaspectratio]{/home/sanjays/et97_scratch2/oldscratch/Ozflux_data_full/ozflux_mirror/www.ozflux.org.au/monitoringsites/titreeeast/images/metaTTE-photo3.jpg}
  \caption*{Photo 3: View of Acacia woodland (foreground) and savanna  (background)}
\end{figure}

\subsection*{Measurement Instrumentation}

\begin{longtable}{p{4.5cm}p{3.5cm}p{3.5cm}p{3.5cm}}
\toprule
\textbf{Instrument Type} & \textbf{Make} & \textbf{Model} & \textbf{Situation} \\
\midrule
\endhead
\bottomrule
\endfoot
Open path  CO2, H2O & Li-COR & Li-7500A & 9.81m \\
Sonic anemometer & Campbell & CSAT3 & 9.81m \\
Air temperature & Vaisala & HMP45C & 9.81m \\
Radiation & Kipp and Zonen & CNR1 & 9.9m \\
Wind speed & R.M. Young & 03101 & 8.28m \\
Wind vane & R.M. Young & 03301 & 8.28m \\
Barometric pressure & Vaisala & CS106 & 2m \\
Rain gauge & Hydrological services & CS700 & 2.5m in savanna \\
Ground heat flux & Middleton & CN3 & -0.08m \\
Soil temperature & Campbell & TCAV & -0.02m – -0.06m  (averaging) \\
Soil moisture – SWC  Reflectometer & Campbell & CS616 & 0m – -0.10m  (averaging) \\
\end{longtable}


\subsection*{Tower Global Metadata}

\begin{longtable}{p{4.5cm}p{10.5cm}}
\toprule
\textbf{Attribute} & \textbf{Value} \\
\midrule
\endhead
\bottomrule
\endfoot
\textbf{PI} & Jamie Cleverly \\
\textbf{Institution} & James Cook University \\
\textbf{Latitude} & -22.283 \\
\textbf{Longitude} & 133.249 \\
\textbf{Elevation} & N/A \\
\textbf{Soil Type} & red kandosol, sandy loam \\
\textbf{Vegetation} & Canopy + three types of understorey (vegetation1-4) \\
\textbf{Temporal Coverage} & 2012-07-18 00:00:00 to 2022-01-17 23:30:00 \\
\textbf{Data Link} & \url{http://data.ozflux.org.au/} \\
\end{longtable}

\newpage

\section{Site Profile: Tumbarumba}

\subsection*{Physical Configuration \& Soil Sensor Depths}

\begin{figure}[H]
\centering
  \includegraphics[width=0.85\textwidth,height=0.45\textheight,keepaspectratio]{/home/sanjays/et97_scratch2/oldscratch/Ozflux_data_full/L6/soil_depth_plots/FluxTower_Diagrams/Tumbarumba/Tumbarumba_diagram.png}
  \caption*{Soil sensor depth profile and tower instrument heights.}
\end{figure}

\newpage

\subsection*{Tower Global Metadata}

\begin{longtable}{p{4.5cm}p{10.5cm}}
\toprule
\textbf{Attribute} & \textbf{Value} \\
\midrule
\endhead
\bottomrule
\endfoot
\textbf{PI} & Will Woodgate (Science); Mark Kitchen (Technical); Markus Loew (Technical) \\
\textbf{Institution} & CSIRO \\
\textbf{Latitude} & -35.6566 \\
\textbf{Longitude} & 148.1516 \\
\textbf{Elevation} & N/A \\
\textbf{Soil Type} & Alfisol \\
\textbf{Vegetation} & wet sclerophyll forest \\
\textbf{Temporal Coverage} & 2002-01-07 16:00:00 to 2022-12-03 18:00:00 \\
\textbf{Data Link} & \url{http://data.ozflux.org.au/} \\
\end{longtable}

\newpage

\section{Site Profile: Warra}

\subsection*{Physical Configuration \& Soil Sensor Depths}

\begin{figure}[H]
\centering
  \includegraphics[width=0.85\textwidth,height=0.45\textheight,keepaspectratio]{/home/sanjays/et97_scratch2/oldscratch/Ozflux_data_full/L6/soil_depth_plots/FluxTower_Diagrams/Warra/Warra_diagram.png}
  \caption*{Soil sensor depth profile and tower instrument heights.}
\end{figure}

\newpage

\subsection*{Site Description}

\textbf{Site Description}

The  Warra flux site is located in the Warra Long-Term Ecological Research (LTER) site  in south western Tasmania (GPS coordinates: -43.0950, 146.6545).

The 15,900 ha Warra LTER  is partly contained within  the Tasmanian Wilderness World Heritage Area (western section of the LTER),  which is managed for conservation, and State forest which is managed for  multipe-uses including wood production.

Timber harvesting on State forest  within the LTER commenced in the 1970's.

Warra  LTER is located between the Huon and Weld Rivers and rises to the Weld Range – a total elevational range of 50-1300 m. The flux site itself is situated on a  small flat at an elevation of 100 m adjacent to the Huon River.

The climate of  Warra classified as temperate with a mild summer and no dry season Mean annual  precipitation is 1700 mm with a relatively uniform seasonal distribution.

Summer  temperatures peak in January (min. 8.4$^\circ$C – max 19.2$^\circ$C) with winter temperatures  reaching their lowest in July (min 2.6$^\circ$C – max 8.4$^\circ$C).

\textit{Eucalyptus obliqua} forests dominate  the vegetation below 650 m where they exist as fire-maintained communities. On  fertile soils these forests attain mature heights in excess of 55m: the tallest \textit{E. obliqua} within the LTER reaches a  height of 90m.

The understorey vegetation progresses from wet sclerophyll  (dominated by \textit{Pomaderris apatala} and \textit{Acacia dealbata}) to rainforest  (dominated by \textit{Nothofagus cunninghamii,  Atherosperma moschatum, Eucryphia lucida and Phyllocladus aspleniifolius})  with increasing time intervals between fire events.

The flux tower is installed  in a stand of tall, mixed-aged \textit{E. obliqua} forest (77 and >250 years-old) with a rainforest understorey and a dense  man-fern (\textit{Dicksonia antarctica})  ground-layer.

The site supports prodigous quantities of coarse woody debris as  is characteristic of these fire-maintained eucalypt forests on fertile sites in  southern Tasmania.

The soil  at the flux site is derived from Permian mudstone and has a gradational profile  with a dark brown organic clayey silt topsoil overlying a yellow brown clay.

The  instruments are mounted at the top of an 80m tall guyed steel lattice tower.  Fluxes of heat, water vapour and carbon dioxide are measured using the  open-path eddy flux technique.

Supplementary measurements above the canopy  include temperature, humidity, windspeed, wind direction, rainfall, incoming  and reflected shortwave radiation and net radiation.

Soil moisture content is  measured using Time Domain reflectometry, while soil heat fluxes and  temperature are also measured.

\subsection*{Purpose \& Scientific Goals}

\textbf{Purpose}

The  purpose of the Warra flux site is:

\begin{itemize}
  \item to study the ecophysiological       processes and rates of C accumulation and decomposition in a mixed-aged,       tall, wet Eucalyptus obliqua forest       originating from past natual wildfires
\end{itemize}

\begin{itemize}
  \item to measure the exchanges of       carbon dioxide, water vapour and energy between the forest and the       atmosphere using micrometeorological techniques
\end{itemize}

\begin{itemize}
  \item to link ecophysiological       processes and rates of C accumulations and decomposition with the biota.
\end{itemize}

\begin{itemize}
  \item to utilize the measurements in       combination with remote sensing data and land surface models to upscale       estimate the net exchanges of carbon and water at regional scale.
\end{itemize}

\subsection*{Site Photographs}

\begin{figure}[H]
\centering
  \includegraphics[width=0.8\textwidth,height=0.4\textheight,keepaspectratio]{/home/sanjays/et97_scratch2/oldscratch/Ozflux_data_full/ozflux_mirror/www.ozflux.org.au/monitoringsites/warra/images/Warra-forest.jpg}
  \caption*{Photo 1: Tall Eucalyptus obliqua wet forest with rainforest understorey at the site of the Warra flux tower}
\end{figure}

\begin{figure}[H]
\centering
  \includegraphics[width=0.8\textwidth,height=0.4\textheight,keepaspectratio]{/home/sanjays/et97_scratch2/oldscratch/Ozflux_data_full/ozflux_mirror/www.ozflux.org.au/monitoringsites/warra/images/TheLoResTower-completed.jpg}
  \caption*{Photo 2: The 80m tall guyed tower houses the flux and meteorological instruments at the Warra flux site}
\end{figure}

\begin{figure}[H]
\centering
  \includegraphics[width=0.8\textwidth,height=0.4\textheight,keepaspectratio]{/home/sanjays/et97_scratch2/oldscratch/Ozflux_data_full/ozflux_mirror/www.ozflux.org.au/monitoringsites/warra/images/IMGP4198.JPG}
  \caption*{Photo 3: Automatic soil chambers measuring trace greenhouse gas fluxes at the Warra Flux site}
\end{figure}

\subsection*{Measurement Instrumentation}

\begin{longtable}{p{4.5cm}p{3.5cm}p{3.5cm}p{3.5cm}}
\toprule
\textbf{Instrument Type} & \textbf{Make} & \textbf{Model} & \textbf{Situation} \\
\midrule
\endhead
\bottomrule
\endfoot
Open path  CO2, H2O & Campbell Scientific & EC150 & 81m \\
3D Sonic anemometer & Campbell Scientific & CSAT3G & 81m \\
Air temperature / relative humidity & Vaisala & HMP155A & 80m \\
Net radiometer & Kipp  and Zonen & CNR4 & 81m \\
Wind speed and direction & Gill & Windsonic4-L 2D sonic & 81m \\
Precipitation & Rimco & RIM-8000 & 80m \\
Soil Heat Flux plate (3 replicates) & Huskeflux & HFP-01 & -0.02m  each \\
Soil  temperature averaging probe                 (3 replicates) & Campbell Scientific & TCAV & -0.02m each \\
Soil temperature thermistor & Campbell Scientific & CS107 & -0.10, -0.20, -0.40, -0.80,  -1.80m \\
Soil  water content reflectometers                 (4x4 depths) & Campbell Scientific & CS616 & -0.10 (x4)                 -0.20 (x2)                  -0.40 (x4)                -0.80 (x2) \\
Power supply & Sealed lead-acid batteries (12V  1050Ah) with 240V chargers linked to petrol generator (Yamaha) &  &  \\
\end{longtable}


\subsection*{Tower Global Metadata}

\begin{longtable}{p{4.5cm}p{10.5cm}}
\toprule
\textbf{Attribute} & \textbf{Value} \\
\midrule
\endhead
\bottomrule
\endfoot
\textbf{PI} & Tim Wardlaw \\
\textbf{Institution} & University of Tasmania \\
\textbf{Latitude} & -43.09502 \\
\textbf{Longitude} & 146.65452 \\
\textbf{Elevation} & N/A \\
\textbf{Soil Type} & Kurosolic Redoxic Hydrosol \\
\textbf{Vegetation} & Eucalyptus obliqua tall forest \\
\textbf{Temporal Coverage} & 2013-03-05 15:00:00 to 2021-09-21 12:00:00 \\
\textbf{Data Link} & \url{http://data.ozflux.org.au/} \\
\end{longtable}

\newpage

\section{Site Profile: Wellington}

\subsection*{Physical Configuration \& Soil Sensor Depths}

\begin{figure}[H]
\centering
  \includegraphics[width=0.85\textwidth,height=0.45\textheight,keepaspectratio]{/home/sanjays/et97_scratch2/oldscratch/Ozflux_data_full/L6/soil_depth_plots/FluxTower_Diagrams/Wellington/Wellington_diagram.png}
  \caption*{Soil sensor depth profile and tower instrument heights.}
\end{figure}

\newpage

\subsection*{Tower Global Metadata}

\begin{longtable}{p{4.5cm}p{10.5cm}}
\toprule
\textbf{Attribute} & \textbf{Value} \\
\midrule
\endhead
\bottomrule
\endfoot
\textbf{PI} & Evan Jensen \\
\textbf{Institution} & University of New South Wales \\
\textbf{Latitude} & -32.5727 \\
\textbf{Longitude} & 148.9843 \\
\textbf{Elevation} & N/A \\
\textbf{Soil Type} &  \\
\textbf{Vegetation} & Pasture \\
\textbf{Temporal Coverage} & 2023-09-27 12:30:00 to 2025-01-17 00:00:00 \\
\textbf{Data Link} & \url{http://data.ozflux.org.au/} \\
\end{longtable}

\newpage

\section{Site Profile: Whroo}

\subsection*{Physical Configuration \& Soil Sensor Depths}

\begin{figure}[H]
\centering
  \includegraphics[width=0.85\textwidth,height=0.45\textheight,keepaspectratio]{/home/sanjays/et97_scratch2/oldscratch/Ozflux_data_full/L6/soil_depth_plots/FluxTower_Diagrams/Whroo/Whroo_diagram.png}
  \caption*{Soil sensor depth profile and tower instrument heights.}
\end{figure}

\newpage

\subsection*{Site Description}

\textbf{Site Description}

The Whroo flux station was located approximately 45km south west of Shepparton, Victoria (GPS coordinates: -36.6732, 145.0294).

The flux tower site was classified as box woodland. The vegetation is dominated by two main Eucalypt species: \textit{Eucalyptus microcarpa} (Grey Box) and \textit{Eucalyptus leucoxylon} (Yellow Gum). Smaller numbers of \textit{Eucalyptus sideroxylon} (Ironbark) and \textit{Acacia pycnantha} (Golden Wattle) are also found on site.

Elevation of the site is close to 165 m and mean annual precipitation from a nearby Bureau of Meteorology site measured 558 mm.

Maximum temperatures ranged from 29.8$^\circ$C (in January) to 12.6$^\circ$C (in July), while minimum temperatures ranged from 14.2$^\circ$C (in February) to 3.2$^\circ$C (in July). Maximum temperatures varied on a seasonal basis by approximately 17.2$^\circ$C and minimum temperatures by 11.0$^\circ$C.

The instrument mast is 36m tall. Heat, water vapour and carbon dioxide measurements are taken using the open-path eddy flux technique.

Temperature, humidity, wind speed, wind direction, rainfall, incoming and reflected shortwave radiation and net radiation were measured above the canopy.

Soil heat fluxes are measured and soil moisture content is gathered using time domain reflectometry.

\subsection*{Purpose \& Scientific Goals}

\textbf{Purpose}

\begin{itemize}
  \item The Whroo site will enable us to assess how revegetation affects  the interaction among carbon dynamics, water regimes (quality, quantity and  frequency) and biodiversity across the landscape (above- and below-ground  terrestrial, and aquatic ecosystems).
  \item We will develop a robust observational and modelling platform  for soil-plant-atmosphere carbon and water fluxes for a range of land-use  practices, including carbon farming, in these landscapes.
  \item We will assess how carbon farming can be modified to increase  biodiversity conservation.
\end{itemize}

\subsection*{Site Photographs}

\begin{figure}[H]
\centering
  \includegraphics[width=0.8\textwidth,height=0.4\textheight,keepaspectratio]{/home/sanjays/et97_scratch2/oldscratch/Ozflux_data_full/ozflux_mirror/www.ozflux.org.au/monitoringsites/whroo/images/Whroo1.jpg}
  \caption*{Photo 1: Above canopy view of the forest}
\end{figure}

\begin{figure}[H]
\centering
  \includegraphics[width=0.8\textwidth,height=0.4\textheight,keepaspectratio]{/home/sanjays/et97_scratch2/oldscratch/Ozflux_data_full/ozflux_mirror/www.ozflux.org.au/monitoringsites/whroo/images/Whroo2.jpg}
  \caption*{Photo 2: The  tower on site}
\end{figure}

\subsection*{Measurement Instrumentation}

\begin{longtable}{p{4.5cm}p{3.5cm}p{3.5cm}p{3.5cm}}
\toprule
\textbf{Instrument Type} & \textbf{Make} & \textbf{Model} & \textbf{Situation} \\
\midrule
\endhead
\bottomrule
\endfoot
Bidirectional  pyranometer / pyrgiometer & Kipp  and Zonen & CNR4 & 36m \\
Combination  cup anemometer / wind vane                (6 replicates) & RM  Young & 03002 & 1m,  2m, 4m, 8m, 16m, 32m \\
Infra-red  gas analyser & LiCOR & LI-7500 & 36m \\
Quantum  sensor & Apogee & SQ-110 &  \\
Soil  heat flux plate                (4 replicates) & Hukseflux & HFP01 & -0.08m \\
Soil  temperature probe                (4 replicates) & Campbell  Scientific & TCAV & All  at -0.08m \\
Solar  regulator & Prostar & PS-30M &  \\
3D Sonic  anemometer & Campbell  Scientific & CSAT3 & 36m \\
Temperature  and Relative Humidity sensor                (6 replicates) & Campbell  Scientific & HMP-45C & 1  each at 1, 2, 4, 8, 16 \& 32m \\
Soil  Moisture Probe                (8 replicates) & Campbell  Scientific & CS-616 & 4  at -0.1m                 \& 1 at each of -0.2, -0.4, -0.8, -1.5m \\
Unidirectional  sunshine pyranometer & Delta & SPN1 & 36m \\
Logger & Campbell  Scientific & CR3000 & 1.5m \\
Multiplexor & Campbell  Scientific & AM25T & 1.5m \\
Rain  gauge & Hydrological Services & TB3 & 36m \\
Profile gas analyser & LiCOR & LI-840 & 32m \\
Profile  pressure line & Vaisala & CS105 & 16m \\
Profile  air pressure & Vaisala & CS106 & 16m \\
Flow  meter & Aalborg & GFM17 & 16m \\
Photosynthetically  Active Radiation (PAR) sensor & Campbell  Scientific & Quantum  LI-190 & 2  at 36m, one facing upward, one facing downward \\
4-channel light sensor & Skye  Instruments & SKR  1850 & 2  at 36m, one facing upward, one facing downward \\
\end{longtable}


\subsection*{Tower Global Metadata}

\begin{longtable}{p{4.5cm}p{10.5cm}}
\toprule
\textbf{Attribute} & \textbf{Value} \\
\midrule
\endhead
\bottomrule
\endfoot
\textbf{PI} & Jason Beringer \\
\textbf{Institution} & University of Western Australia \\
\textbf{Latitude} & -36.6732 \\
\textbf{Longitude} & 145.0294 \\
\textbf{Elevation} & N/A \\
\textbf{Soil Type} & ?? \\
\textbf{Vegetation} & Dry schlerophyll woodland \\
\textbf{Temporal Coverage} & 2011-12-01 13:30:00 to 2023-12-31 23:30:00 \\
\textbf{Data Link} & \url{http://data.ozflux.org.au/} \\
\end{longtable}

\newpage

\section{Site Profile: Yanco}

\subsection*{Physical Configuration \& Soil Sensor Depths}

\begin{figure}[H]
\centering
  \includegraphics[width=0.85\textwidth,height=0.45\textheight,keepaspectratio]{/home/sanjays/et97_scratch2/oldscratch/Ozflux_data_full/L6/soil_depth_plots/FluxTower_Diagrams/Yanco/Yanco_diagram.png}
  \caption*{Soil sensor depth profile and tower instrument heights.}
\end{figure}

\newpage

\subsection*{Site Description}

\textbf{Site Description}

The flux tower monitoring system is located within the Yanco region of New South Wales, Australia (GPS coordinates: -34.9893, 146.2907). The Yanco area of the flux tower site is located in the western plains of the Murrumbidgee Catchment and is within a wider research area (60 x 60 km) that supports a network of OzNet stations, which have been in operation since late 2001 onwards.

This is a topographically flat area, primarily comprised of the following soil types: sandy loams, scattered clays, red brown earths, transitional red brown earth, sands over clay and deep sands. Stream valleys and layered soil and sedimentary materials are found across the landscape.

The tower on site extends to 20m, however flux measurements are recorded from slightly lower than this (see measurements table for more details).

Mean annual precipitation from a nearby Bureau of Meteorology site measured 465 mm.

Maximum temperatures ranged from 37.4$^\circ$C (in January) to 16.6$^\circ$C (in July), while minimum temperatures ranged from 29.0$^\circ$C (in January) to 11.8$^\circ$C (in July). Maximum temperatures varied on a seasonal basis by approximately 20.8$^\circ$C and minimum temperatures by 17.2$^\circ$C.

\subsection*{Purpose \& Scientific Goals}

\textbf{Purpose}

\begin{itemize}
  \item To provide additional       validation data for use in land surface data assimilation. These data systems are important for land surface       models, which rely on atmosphere-surface observations.
  \item For validation of       the Advanced Microwave Scanning Radiometer-2 (AMSR2) soil moisture product       and resulting soul moisture and flux estimates.
\end{itemize}

\begin{itemize}
  \item For crop modelling, runoff, flood forecasting, Land  Surface Model forcing data (i.e. input of AMSR2 soil moisture into the Joint UK  Land Environment Simulator (JULES) model).
\end{itemize}

\subsection*{Site Photographs}

\begin{figure}[H]
\centering
  \includegraphics[width=0.8\textwidth,height=0.4\textheight,keepaspectratio]{/home/sanjays/et97_scratch2/oldscratch/Ozflux_data_full/ozflux_mirror/www.ozflux.org.au/monitoringsites/yanco/images/Jaxa1.JPG}
  \caption*{Photo 1: The  tower on site}
\end{figure}

\begin{figure}[H]
\centering
  \includegraphics[width=0.8\textwidth,height=0.4\textheight,keepaspectratio]{/home/sanjays/et97_scratch2/oldscratch/Ozflux_data_full/ozflux_mirror/www.ozflux.org.au/monitoringsites/yanco/images/Jaxa2.JPG}
  \caption*{Photo 2: Map  of the flux tower and other research sites in the Yanco region}
\end{figure}

\subsection*{Measurement Instrumentation}

\begin{longtable}{p{4.5cm}p{3.5cm}p{3.5cm}p{3.5cm}}
\toprule
\textbf{Instrument Type} & \textbf{Make} & \textbf{Model} & \textbf{Situation} \\
\midrule
\endhead
\bottomrule
\endfoot
Pyranometer  - with defrosting fan & Kipp and Zonen & CMP21 & 2m \\
Pyrgeometer   - with defrosting fan & Kipp and Zonen & CGR4 & 2m \\
PAR quantum sensor & Kipp and Zonen & PQS1 & 2m \\
Infrared radiation  thermometer & Apogee & SI-111 & 0.5m \\
3D sonic anemometer & Campbell Scientific & CSAT3 & 8m \\
H20/CO2 analyzer & LiCOR & 7500 & 8m \\
Wind sensor & Campbell Scientific & 034B & 4  at 2, 4, 8, 16m \\
Temperature and humidity  probe with ventilation unit & Campbell Scientific & HMP-45C,  PCV-04 & 4  at 2, 4, 8, 16m \\
Rain gauge & TKF-1-UD & TKF-1-UD & 0.5m \\
Barometer & Vaisala & PTB210 & 8m \\
Data logger & Campbell Scientific & CR1000-4M, AM16, NL115 & 4m \\
Transformer DC12V-DV24V &  & EX-375L2 & 4m \\
DC power supply &  & CS616 & 0.5m \\
Automatic station for soil  hydrology components &  & ASSH-T & 0.5m \\
Solar panel & Campbell Scientific & C-MSX10 & 1m \\
Multi-weather sensor & Vaisala & Weather transmitter, WTX520 & 6m \\
Soil temperature sensor & WST-110 & WST-110 & 5 probes  installed at 0.03, 0.1, 0.15, 0.45 \& 0.75m \\
Soil moisture sensor & IMKO & TRIME-PICO32 & 5 probes  installed at 0.03, 0.1, 0.15, 0.45 \& 0.75m \\
Soil heat flux sensor & Hukseflux & HF-HFP01 & 5  plates, 3 at 0.03m,                2 at 0.07m \\
\end{longtable}


\subsection*{Tower Global Metadata}

\begin{longtable}{p{4.5cm}p{10.5cm}}
\toprule
\textbf{Attribute} & \textbf{Value} \\
\midrule
\endhead
\bottomrule
\endfoot
\textbf{PI} & Jason Beringer \\
\textbf{Institution} & Monash University and The University of Western Australia \\
\textbf{Latitude} & -34.9878 \\
\textbf{Longitude} & 146.2908 \\
\textbf{Elevation} & N/A \\
\textbf{Soil Type} & sandy loams, scattered clays, red brown earths, deep sands \\
\textbf{Vegetation} & Pasture \\
\textbf{Temporal Coverage} & 2012-10-24 14:00:00 to 2025-01-28 03:00:00 \\
\textbf{Data Link} & \url{http://data.ozflux.org.au/} \\
\end{longtable}

\newpage

\end{document}
